\documentclass[11pt,a4paper]{article}
\usepackage[T1]{fontenc}
\usepackage[utf8]{inputenc}
\usepackage{amsmath,amssymb,amsthm,mathtools}
\usepackage{booktabs}
\usepackage{microtype}
\usepackage[margin=1.05in]{geometry}
\usepackage{xcolor}
\usepackage[colorlinks=true,linkcolor=blue!50!black,citecolor=blue!50!black,urlcolor=blue!50!black]{hyperref}
\usepackage{cleveref}

\newtheorem{theorem}{Theorem}[section]
\newtheorem{lemma}[theorem]{Lemma}
\newtheorem{proposition}[theorem]{Proposition}
\newtheorem{corollary}[theorem]{Corollary}
\theoremstyle{definition}
\newtheorem{definition}[theorem]{Definition}
\newtheorem{example}[theorem]{Example}
\theoremstyle{remark}
\newtheorem{remark}[theorem]{Remark}
\crefname{theorem}{Theorem}{Theorems}
\crefname{lemma}{Lemma}{Lemmas}
\crefname{proposition}{Proposition}{Propositions}
\crefname{corollary}{Corollary}{Corollaries}
\crefname{definition}{Definition}{Definitions}
\crefname{remark}{Remark}{Remarks}
\crefname{equation}{}{}
\crefformat{equation}{(#2#1#3)}

\newcommand{\R}{\mathbb{R}}
\newcommand{\E}{\mathbb{E}}
\newcommand{\PP}{\mathbb{P}}
\newcommand{\one}{\mathbf{1}}
\newcommand{\norm}[1]{\left\lVert #1\right\rVert}
\newcommand{\abs}[1]{\left\lvert #1\right\rvert}
\newcommand{\op}[1]{\norm{#1}_{2\to 2}}
\DeclareMathOperator{\Lip}{Lip}
\DeclareMathOperator{\im}{im}
\DeclareMathOperator{\Var}{Var}
\DeclareMathOperator{\diag}{diag}
\DeclareMathOperator{\rank}{rank}
\newcommand{\dd}{\,\mathrm{d}}
\newcommand{\eqd}{\overset{\mathrm{d}}{=}}
\newcommand{\gauge}{\mathcal{G}}
\newcommand{\net}{\mathcal{N}}

\title{Global Stability of Deep Gaussian ReLU Networks}
\author{Yitzchak Shmalo\\[3pt]
\small Einstein Institute of Mathematics, The Hebrew University of Jerusalem, Jerusalem, Israel\\
\small \texttt{yitzchak.shmalo@gmail.com}}
\date{}

\begin{document}
\maketitle

\begin{abstract}
A bias-free ReLU network of width $n$ and depth $L$ with independent $\mathcal N(0,\alpha/n)$ weights has a global Lipschitz constant $K_{n,L}(\alpha)$, the worst stretching over all pairs of inputs for one draw of the weights. We prove, for every $n$, every $L\ge2$ and every $\theta>0$,
\[
\PP\bigl(K_{n,L}(\alpha)>1\bigr)\le2\exp\!\bigl(n\,B_L(\alpha,\theta)\bigr),\qquad B_L=L\,H(1/L)+\log5+\theta\log4+L\,g_\alpha(\theta),
\]
with $H$ the binary entropy and $g_\alpha$ elementary, and for $\alpha>2$ the converse bound $\PP(K_{n,L}(\alpha)\le1)\le5\alpha^2L/(n(\alpha-2)^2)$. The fixed-depth threshold $\alpha_L$ therefore tends to $2$, with $\alpha_L\ge2\exp[-(\sqrt{10}+o(1))\sqrt{\log L/L}]$. The proof turns on an exact identity for realizable activation patterns: averaging over the $2^{nL}$ sign gauges, which preserve the law of the weights, turns the realizability event into a deterministic orthant count divided by $2^{nL}$.

The second result ties the layers: for one matrix $W$ applied repeatedly, $x_{t+1}=\phi(Wx_t)$, and any fixed horizon $T$, $\max_{t\le T}|2^t\norm{x_t}^2-1|\le\delta$ with probability at least $1-Ce^{-cN}$ for every deterministic unit start. For a typical matrix, the bound holds outside a set of exponentially small measure under any prescribed input distribution independent of the matrix. Here the dependence of the trajectory on the matrix is handled by conditioning, made exact by a Gaussian innovation lemma for adaptive orthogonal queries. The two results share their object and their obstacle and resolve it in opposite ways; the paper explains why each method fits its question, and uses the second to sharpen the expansive side of the first to an exponential rate. Selected identities, bounds and norm statistics are checked numerically at small width.
\end{abstract}

\tableofcontents

\section{The object, the question, and the answer}
\label{sec:intro}

\subsection{The object}

Fix integers $n\ge 1$ and $L\ge 1$ and a number $\alpha>0$. Let $W_1,\dots,W_L$ be independent $n\times n$ random matrices, every entry of every matrix independent with law $\mathcal N(0,\alpha/n)$. Write $\phi(t)=\max\{t,0\}$ for the scalar ReLU and apply it to a vector coordinate by coordinate. Starting from an input $s\in\R^n$ set
\begin{equation}\label{eq:forward}
x_0(s)=s,\qquad h_\ell(s)=W_\ell x_{\ell-1}(s),\qquad x_\ell(s)=\phi(h_\ell(s)),\qquad 1\le \ell\le L,
\end{equation}
and let $F^{(\alpha)}_{n,L}(s)=x_L(s)$. There are no biases, no training and no data: the network is $L$ random matrices and the operation of clipping negative coordinates to zero. The factor $1/n$ in the variance is the usual width normalization; it keeps a coordinate of $W_\ell x$ at variance $(\alpha/n)\norm{x}^2$, so that $\norm{W_\ell x}^2$ has expectation $\alpha\norm{x}^2$ rather than $n\alpha\norm x^2$. The parameter $\alpha$ is the dial of the problem.

\subsection{The question}

A function $F:\R^n\to\R^n$ stretches or shrinks distances. Its global Lipschitz constant
\begin{equation}\label{eq:Kdef}
K_{n,L}(\alpha)=\Lip_2\bigl(F^{(\alpha)}_{n,L}\bigr)=\sup_{s\ne t}\frac{\norm{F^{(\alpha)}_{n,L}(s)-F^{(\alpha)}_{n,L}(t)}}{\norm{s-t}}
\end{equation}
is the worst stretching factor over every pair of inputs. It is not an average and not a typical value; one realization of the weights must control every pair at once. When $K_{n,L}(\alpha)\le 1$ the network never increases a distance and we call it non-expansive. The question is: for which $\alpha$ is the random network non-expansive, and how does the answer depend on the width and the depth?

For a fixed depth $L$, define the fixed-depth threshold
\begin{equation}\label{eq:alphaL}
\alpha_L=\sup\bigl\{\alpha>0:\ \PP\bigl(K_{n,L}(\alpha)\le 1\bigr)\to 1\ \text{as } n\to\infty\bigr\}.
\end{equation}
Membership in the set is itself a statement about a limit in the width, with $\alpha$ and $L$ held fixed.

\subsection{The answer, at finite width}

The proof below produces, for every width and depth, an explicit bound on the probability that the network is expansive. Two elementary functions carry it. The binary entropy is
\[
H(u)=-u\log u-(1-u)\log(1-u),\qquad 0\le u\le 1,\qquad 0\log 0=0,
\]
and for $0<\theta<1/4$,
\begin{equation}\label{eq:hdef}
h_\alpha(\theta)=\theta\log\frac{\alpha}{2e}+\log\Bigl(\frac{1+(1-4\theta)^{-1/2}}{2}\Bigr).
\end{equation}
The entropy will measure how many activation patterns a random network can realize; $h_\alpha$ will measure the high moments of one layer's radial factor.

\begin{theorem}[finite width, finite depth]\label{thm:finite}
For every $n\ge 1$, every $L\ge 2$, every $\alpha>0$ and every $\theta\in(0,1/4)$,
\begin{equation}\label{eq:main-finite}
\begin{gathered}
\PP\bigl(K_{n,L}(\alpha)>1\bigr)\;\le\;\E\bigl[K_{n,L}(\alpha)^{2\theta n}\bigr]\;\le\;2\exp\bigl(n\,B_L(\alpha,\theta)\bigr),\\
B_L(\alpha,\theta)=L\,H(1/L)+\log 5+\theta\log 4+L\,h_\alpha(\theta).
\end{gathered}
\end{equation}
For $L=1$ the same holds with $L H(1/L)$ replaced by $\log 2$.
\end{theorem}

The four terms of the budget $B_L$ have four separate origins, listed in \cref{sec:budget}: the entropy term counts realizable activation patterns, $\log 5$ pays for a net of the sphere, $\theta\log 4$ for replacing an operator norm by a maximum over that net, and $L\,h_\alpha(\theta)$ bounds the radial moment exponent over $L$ layers. The first three do not depend on $\alpha$; the last is the only one that carries a factor $L$ with a sign that can be negative. When $\alpha<2$ the function $h_\alpha$ is negative on an interval $(0,\theta_0)$, because its derivative at zero is $\log(\alpha/2)<0$, so the last term is a negative multiple of $L$ while the positive terms grow like $\log L$. Beyond a computable depth the budget is negative.

\begin{corollary}[fixed depth, growing width]\label{cor:width}
Let $\alpha<2$ and $L\ge 2$, and suppose $c_L(\alpha):=\max_{0<\theta<1/4}\bigl[-B_L(\alpha,\theta)\bigr]>0$. Then for every width $n$,
\[
\PP\bigl(K_{n,L}(\alpha)>1\bigr)\le 2e^{-c_L(\alpha)\,n},
\]
so the network is non-expansive with probability at least $1-2e^{-c_L(\alpha)n}$. In particular $\PP(K_{n,L}(\alpha)\le 1)\to 1$ as $n\to\infty$, with an exponential rate.
\end{corollary}

At a given finite width the probability is not one; the bound says how far from one it can be. For instance at $\alpha=1$ and $L=1000$ the rate is $c_{1000}(1)=25.03$, and already at width $n=10$ the bound reads $2e^{-250}$. \Cref{sec:numerics} tabulates these numbers.

\begin{theorem}[the other side, at finite width]\label{thm:above2}
For every $n\ge 1$, $L\ge 1$ and $\alpha>2$,
\[
\PP\bigl(K_{n,L}(\alpha)\le 1\bigr)\le\frac{5\alpha^2 L}{n(\alpha-2)^2}.
\]
\end{theorem}

The two finite statements combine into the theorem the paper is named for.

\begin{theorem}[unconditional global threshold]\label{thm:main}
For the Gaussian bias-free ReLU network above:
\begin{enumerate}
\item[(i)] For every $\alpha<2$ there is a depth $L_0(\alpha)$ such that for every fixed $L\ge L_0$ there is $c_{\alpha,L}>0$ with $\PP(K_{n,L}(\alpha)>1)\le 2e^{-c_{\alpha,L}n}$ for all $n$; in particular $\PP(K_{n,L}(\alpha)\le 1)\to 1$.
\item[(ii)] For every fixed $L\ge 1$ and every $\alpha>2$, $\PP(K_{n,L}(\alpha)>1)\to 1$ as $n\to\infty$.
\item[(iii)] Consequently $\lim_{L\to\infty}\alpha_L=2$.
\item[(iv)] More precisely, $\alpha_L\ge 2\exp\bigl[-(\sqrt{10}+o(1))\sqrt{\log L/L}\bigr]$ as $L\to\infty$.
\end{enumerate}
\end{theorem}

Part (i) is \cref{cor:width} together with the observation that the budget turns negative at large depth; the depth $L_0(\alpha)$ is the smallest $L$ at which $c_L(\alpha)>0$, and it is computable (\cref{sec:numerics}). Part (ii) is \cref{thm:above2}. Part (iii) follows from (i) and (ii). Part (iv) comes from expanding $h_\alpha$ near $\alpha=2$, where $h_{2e^{-\delta}}(\theta)=-\delta\theta+\tfrac52\theta^2+O(\theta^3)$; optimizing $\theta$ and balancing the negative term $-L\delta^2/10$ against the entropy $\log L$ gives $\delta\approx\sqrt{10\log L/L}$.

\subsection{The order of limits}

The order is part of the statement. First a depth $L$ is fixed; then the width $n$ tends to infinity, which defines $\alpha_L$; only then does $L$ tend to infinity. Nothing is claimed about a joint regime $L=L(n)$. \Cref{thm:finite} is nevertheless a finite statement in both variables, and one may read it in any regime one likes, provided the budget $B_L$ is evaluated there. For example, at fixed $\alpha<2$ the bound is useless below the depth $L_0(\alpha)$, whatever the width, which is a limitation of this certificate. There is also a genuine depth dependence: at depth one the threshold is $\alpha_1=1/4$, so a depth-one network at $\alpha=1$ is expansive with probability tending to one although $1<2$.

\subsection{Why the number two}

For a fixed nonzero $x$, each coordinate of $W_\ell x$ is Gaussian with variance $(\alpha/n)\norm x^2$, and ReLU keeps, in expectation, half of the squared magnitude:
\[
\E\bigl[\norm{\phi(W_\ell x)}^2\ \big|\ x\bigr]=\frac{\alpha}{2}\norm x^2.
\]
A forward pass multiplies a typical squared norm by $\alpha/2$ per layer, and $\alpha=2$ is where that factor is one. The content of the theorem is that the same number governs the global Lipschitz constant in the deep limit, although the supremum in \cref{eq:Kdef} may hunt for rare activation regions and rare directions, and although the region an input lands in is decided by the very matrices whose stretching is being measured.

\subsection{The second result: one matrix, many steps}

The network above draws a fresh matrix at every layer. Tie the layers instead, one matrix $W$ with $\mathcal N(0,1/N)$ entries applied again and again, $x_{t+1}=\phi(Wx_t)$ from a unit vector $x_0$; in the notation above this is $\alpha=1$ and $n=N$. At $\alpha=1$ a layer halves the expected squared norm, and the second theorem (\cref{sec:horizon}) says that this halving is exact to any tolerance, uniformly over a fixed number of steps, with exponentially small failure probability: for every $T$ and $\delta$ there are $C,c$ depending only on them with
\[
\PP\Bigl(\max_{0\le t\le T}\bigl|2^t\norm{x_t}^2-1\bigr|>\delta\Bigr)\le Ce^{-cN}
\]
for every deterministic $x_0$; and for one typical matrix the set of starting points that fail has exponentially small measure under any measure independent of $W$. The obstacle is the same as before: $x_1$ is a function of $W$, so $Wx_1$ is not a fresh Gaussian vector. The resolution is the opposite of the gauge argument. The trajectory spans a subspace of dimension at most $T+1$; a query $Wu$ reveals only the coefficients of the rows of $W$ in the direction $u$; and if the directions queried are orthogonal and each is chosen from what has already been revealed, the answers are independent Gaussian vectors, whatever the rule that chose them. Gram--Schmidt on the trajectory produces such directions, the dynamics is rewritten in terms of $T+1$ Gaussian innovations, and a law of large numbers uniform over a ball in $\R^{T+1}$ closes the argument. \Cref{sec:together} sets the two results side by side: what they share, where they differ, why each method fits its question, and one place where the second method improves the first theorem.

\subsection{What is new here}

The first theorem's proof turns on the exact orbit identity of \cref{sec:gauge}: the probability that a formal activation pattern is realized, weighted by any gauge-invariant functional, equals the expected number of orthants met by a random subspace divided by $2^{nL}$. Nothing is conditioned on and no independence is assumed; the identity is exact and the only inequality is a deterministic orthant count. Written at finite width, the argument gives an explicit failure bound for every $n$ and $L$, with a budget whose four terms have four named origins, and a computable sufficient finite-depth threshold; the limit theorem and a lower bound on the threshold are read off from it. The second theorem's proof turns on the innovation lemma of \cref{sec:horizon}, which makes conditioning on an adaptively chosen trajectory exact. Placing the two together shows that the same entanglement admits two different exact treatments, that the choice between them is decided by what is quantified over, one input or all inputs, and that the second treatment sharpens the expansive side of the first theorem.


\subsection{Related work and the scope of the contribution}
\label{sec:related}

Global Lipschitz bounds for random ReLU networks already have a substantial literature. Geuchen, St\"oger, Telaar and Voigtlaender~\cite{geuchen2025} study scalar-output maps $\Phi:\R^d\to\R$ with hidden width $N$, Gaussian weights and independent biases. Their Theorems~3.3 and~3.4 give upper and lower bounds for deep networks: the upper bound assumes $N>d+2$, and the lower bound assumes $N\gtrsim dL^2$. Their hidden-weight variance is $2/N$ and their final scalar readout has variance one; normalization therefore matters when comparing numerical Lipschitz constants. Our map has input, hidden and output dimension all equal to $n$, variance $\alpha/n$, and ReLU at the output. We control its Euclidean operator norm globally and obtain a sufficient contraction certificate at every finite width. The cited deep-network theorems cannot simply be specialized to this square regime by setting $d=N$.

Dirksen, Finke, Geuchen, St\"oger and Voigtlaender~\cite{dirksen2025} sharpen the scalar-output theory for all $\ell^p$ input metrics. In the zero-bias case, their Corollary~4.2 gives, for $p=2$, an upper bound of order $\sqrt{d\log(N/d)}$ under a width condition of order $N\gtrsim dL^3$, up to logarithmic factors; their Corollary~3.6 gives a lower bound of order $\sqrt d$ under its hypotheses. Their global upper bound combines gradient estimates with control of nearby activation changes using random hyperplane tessellations. Our argument instead averages strict realizability over sign gauges, retaining its dependence on the Jacobian functional, before applying a net to a fixed formal linear map. The different dimensions, readout and width hypotheses prevent a direct numerical comparison or a claim that one theorem improves the other.

For tied weights, Yang's Tensor Programs~I~\cite{yang2019}, in particular Theorem~5.4 and its recurrent-network examples, supplies a broad framework for Gaussian limits and convergence of empirical coordinate averages in fixed programs. Reusing Gaussian matrices and deriving deterministic norm recursions are therefore not new in themselves. Our finite-horizon statement gives a quantitative failure bound $Ce^{-cN}$ for this particular ReLU iteration, uniformly over a fixed time horizon and uniformly in the choice of deterministic unit start. A Fubini argument then bounds the measure of bad starts for a typical matrix. Neither this conclusion nor the cited fixed-program limit supplies uniform control over every input of one matrix or all-time convergence.

The contribution claimed here is the weighted sign-gauge identity and the finite-width budget it yields for square vector-output networks, together with a self-contained quantitative finite-horizon argument and an explicit comparison of their quantifiers. The rate in \cref{thm:main}(iv) is a certified lower bound on $\alpha_L$, not a matching asymptotic rate for the true threshold.

\subsection{How the proof goes}

The difficulty is an entanglement. The activation pattern of an input is a function of the weights; the worst input is chosen after the weights are seen; and the Jacobian on the region of that input is built from the same weights. Conditioning on a realized pattern would tilt the law of the weights in a way that is hard to control. The proof never conditions. Its steps:
\begin{enumerate}
\item Replace the global supremum by the maximum operator norm of the linear maps on the nonempty activation regions (\cref{sec:regions}).
\item Fix a \emph{formal} mask sequence $D$: a sequence of diagonal $0/1$ matrices written down without asking whether any input produces it. Realizability of $D$ is the event that an $n$-dimensional subspace of $\R^{nL}$ meets one specified open orthant.
\item Average that event over a group of $2^{nL}$ sign changes. The sign changes preserve the joint law of the weights and the norm of every formal Jacobian, while they move the target orthant through all $2^{nL}$ orthants exactly once. The average of the realizability indicator is therefore a deterministic orthant count divided by $2^{nL}$, and an exact identity results (\cref{sec:gauge}). This identity is the centre of the argument.
\item Bound the orthant count by a classical result: an at most $n$-dimensional subspace of $\R^{nL}$ meets at most $C_{n,L}=2\sum_{j<n}\binom{nL-1}{j}\le 2e^{nLH(1/L)}$ open orthants.
\item Bound the operator norm of a formal Jacobian through a $\tfrac12$-net of the sphere and exact Gaussian radial moments, and reduce everything to one scalar random variable $X_n^{(\alpha)}=\frac{\alpha}{n}\sum_i B_iZ_i^2$ (\cref{sec:moments}).
\item Take a moment of order $2\theta n$ and bound it with an explicit finite-$n$ inequality; the four exponential rates assemble into $B_L$ (\cref{sec:theorems}).
\end{enumerate}
Every estimate is finite in $n$ and $L$. The only place a limit enters is the definition of $\alpha_L$ itself, and the asymptotic form (iv), where an expansion in $\theta$ is convenient. \Cref{sec:numerics} checks the identity of step 3 and the bounds of steps 4--6 against brute-force computation at small width, and \cref{sec:sharp} computes the exact scalar exponent $g_\alpha$. The finite upper bound $h_\alpha$ agrees with $g_\alpha$ through second order at $\theta=0$, not for general $\theta$: $g_\alpha(\theta)=h_\alpha(\theta)+O(\theta^3)$. Thus the derivative $\log(\alpha/2)$ and quadratic coefficient $\tfrac52$ are shared, while the finite theorem uses the upper bound $h_\alpha$. For example, at $\alpha=2$ and $\theta=0.1$, $h_2=0.0358388$ whereas $g_2=0.0210736$.

\section{Tools}
\label{sec:tools}

Everything in this section is standard. It is collected so that the later sections can refer to one place and so that the paper reads without outside references.

\subsection{Norms}

For $x\in\R^n$, $\norm x=\bigl(\sum_i x_i^2\bigr)^{1/2}$ and $\langle x,y\rangle=x^Ty$; the Cauchy--Schwarz inequality $\abs{\langle x,y\rangle}\le\norm x\norm y$ says a cosine is at most one. The unit sphere is $S^{n-1}=\{x:\norm x=1\}$. For a linear map $T:\R^n\to\R^m$ the operator norm is
\[
\op{T}=\sup_{x\ne 0}\frac{\norm{Tx}}{\norm x}=\sup_{\norm x=1}\norm{Tx},
\]
the two forms agreeing because $x=\norm x\cdot x/\norm x$. Thus $\norm{Tx}\le\op T\norm x$ for every $x$. The supremum is attained on the compact sphere; a unit vector $v$ with $\norm{Tv}=\op T$ is a right singular vector for the largest singular value. If $Q$ is orthogonal, $Q^TQ=I$, then $\norm{Qx}^2=x^TQ^TQx=\norm x^2$, so $\op{QT}=\op T=\op{TQ}$. A diagonal matrix with entries $\pm1$ is orthogonal; this is the one fact about orthogonal matrices the proof uses.

\subsection{ReLU}

$\phi(t)=\max\{t,0\}$, applied coordinatewise to vectors.

\begin{lemma}\label{lem:relu}
For all real $a,b$ and every $c\ge 0$: $\abs{\phi(a)-\phi(b)}\le\abs{a-b}$, $\phi(ca)=c\phi(a)$, and $\phi(a)=\one_{\{a>0\}}a$ whenever $a\ne 0$. Consequently $\norm{\phi(x)-\phi(y)}\le\norm{x-y}$ for vectors.
\end{lemma}

\begin{proof}
If $a,b\ge 0$ the first inequality is an equality; if $a,b\le 0$ both values are zero; if $a\ge 0\ge b$ then $\abs{\phi(a)-\phi(b)}=a\le a-b$. Homogeneity holds because multiplying by $c\ge 0$ does not change the sign of $a$, and the indicator formula is the definition away from zero. For vectors, square the scalar inequality coordinatewise and sum.
\end{proof}

Homogeneity is what makes a bias-free network positively homogeneous: from \cref{eq:forward}, $x_\ell(cs)=cx_\ell(s)$ for $c\ge 0$ by induction on $\ell$, so
\begin{equation}\label{eq:homog}
F^{(\alpha)}_{n,L}(cs)=cF^{(\alpha)}_{n,L}(s),\qquad c\ge 0,\qquad\text{and}\qquad F^{(\alpha)}_{n,L}(0)=0.
\end{equation}

\subsection{Lipschitz constants and Jacobians}

$F:\R^n\to\R^m$ is globally Lipschitz if $\norm{F(s)-F(t)}\le K\norm{s-t}$ for all $s,t$; the least such $K$ is $\Lip_2(F)$. If $F$ is differentiable at $s$ with Jacobian $DF(s)$, then for a unit vector $v$,
\[
\norm{DF(s)v}=\lim_{h\to 0}\frac{\norm{F(s+hv)-F(s)}}{\abs h}\le\Lip_2(F),
\]
so $\op{DF(s)}\le\Lip_2(F)$ at every point of differentiability. For a continuous map that is linear on each of finitely many polyhedral pieces the reverse also holds, and equality is exactly what the network needs; that is \cref{lem:pl} below.

The constant $K_{n,L}(\alpha)$ of \cref{eq:Kdef} is finite: ReLU is $1$-Lipschitz, so
\[
\norm{x_\ell(s)-x_\ell(t)}\le\op{W_\ell}\norm{x_{\ell-1}(s)-x_{\ell-1}(t)},
\]
and iterating gives $K_{n,L}(\alpha)\le\prod_\ell\op{W_\ell}<\infty$. It is measurable: by continuity the supremum in \cref{eq:Kdef} may be restricted to rational pairs, and each ratio is a Borel function of the weights.

\subsection{Three probability inequalities}

\begin{lemma}\label{lem:prob}
(a) Union bound: $\PP(\bigcup_j A_j)\le\sum_j\PP(A_j)$. (b) Markov: if $Y\ge 0$ and $a>0$ then $\PP(Y>a)\le\E Y/a$; in particular $\PP(Y>1)\le\E Y$. (c) Chebyshev: if $Y$ has finite variance then $\PP(\abs{Y-\E Y}>\varepsilon)\le\Var(Y)/\varepsilon^2$.
\end{lemma}

\begin{proof}
(a) $\one_{\cup A_j}\le\sum\one_{A_j}$ pointwise; take expectations. (b) $Y\ge a\one_{\{Y>a\}}$ pointwise; take expectations and divide by $a$. (c) Apply (b) to $(Y-\E Y)^2$ with threshold $\varepsilon^2$.
\end{proof}

\subsection{Gaussian facts}

$Z\sim\mathcal N(0,1)$ has density $(2\pi)^{-1/2}e^{-z^2/2}$, $\E Z^2=1$, $\E Z^4=3$ (two integrations by parts), and for $t<1/2$
\begin{equation}\label{eq:mgf}
\E e^{tZ^2}=\frac{1}{\sqrt{2\pi}}\int_\R e^{-(1-2t)z^2/2}\dd z=(1-2t)^{-1/2}.
\end{equation}
If $G\in\R^m$ has independent $\mathcal N(0,1)$ coordinates, $\norm G^2$ has the chi-square law with $m$ degrees of freedom, written $\chi^2_m$; we set $\chi^2_0=0$. If $Q$ is a fixed orthogonal matrix then $QG$ is again centred Gaussian with covariance $QQ^T=I$, so $QG\eqd G$; in particular multiplying rows or columns of a Gaussian matrix by fixed signs does not change its law. A centred Gaussian vector is determined by its covariance, so jointly Gaussian variables with zero covariance are independent.

\subsection{A bound on partial binomial sums}

\begin{lemma}\label{lem:entropy}
For integers $N\ge1$ and $0\le k\le N/2$,
\[
\sum_{j=0}^{k}\binom Nj\le\exp\bigl(N\,H(k/N)\bigr).
\]
\end{lemma}

\begin{proof}
For $k=0$ both sides are $1$. Otherwise put $q=k/N\in(0,1/2]$. Expanding $1=(q+(1-q))^N$ and keeping only the terms $j\le k$,
\[
1\ge\sum_{j\le k}\binom Nj q^j(1-q)^{N-j}\ge\sum_{j\le k}\binom Nj q^k(1-q)^{N-k},
\]
because $q\le 1-q$ makes $q^j(1-q)^{N-j}$ non-increasing in $j$. Hence $\sum_{j\le k}\binom Nj\le q^{-k}(1-q)^{-(N-k)}=\exp(NH(q))$.
\end{proof}

This is the only counting estimate the finite theorem needs. Stirling's formula appears in \cref{sec:sharp}, where the exact exponent is computed, and nowhere else.

\subsection{One deterministic inequality}

\begin{lemma}\label{lem:xp}
For every $x\ge 0$, $p>0$ and $t>0$, $x^p\le(p/(et))^pe^{tx}$.
\end{lemma}

\begin{proof}
For $x>0$, $\log(x^pe^{-tx})=p\log x-tx$ has derivative $p/x-t$, positive before $x=p/t$ and negative after, so the maximum of $x^pe^{-tx}$ is $(p/t)^pe^{-p}$. Multiply by $e^{tx}$.
\end{proof}

The inequality is tight at exactly one point, $x=p/t$; it is what converts a moment of order $p$ into a moment-generating function, at the price of the factor $(p/(et))^p$.

\section{Regions: from a global supremum to formal masks}
\label{sec:regions}

\subsection{Activation masks and the Jacobian}

At an input $s$ for which no coordinate of any preactivation $h_\ell(s)$ is zero, define the activation masks
\[
D_\ell(s)=\diag\bigl(\one_{\{h_{\ell 1}(s)>0\}},\dots,\one_{\{h_{\ell n}(s)>0\}}\bigr),\qquad 1\le\ell\le L,
\]
diagonal $0/1$ matrices recording which coordinates the ReLU keeps. Away from zero the derivative of the scalar ReLU is $1$ on the positive half-line and $0$ on the negative one, so the chain rule gives
\begin{equation}\label{eq:jac}
J(s)=D_L(s)\,W_L\,D_{L-1}(s)\,W_{L-1}\cdots D_1(s)\,W_1 .
\end{equation}
The network is linear in a neighbourhood of such an $s$, with matrix $J(s)$. The difficulty named in the introduction is visible in \cref{eq:jac}: the masks are functions of the weights, and the input that produces the worst mask is chosen after the weights are drawn.

\subsection{Formal masks}

The step that breaks the entanglement is to stop insisting that a mask come from an input.

\begin{definition}[formal mask sequence]
A formal mask sequence is a tuple $D=(D_1,\dots,D_L)$ of $n\times n$ diagonal matrices with diagonal entries in $\{0,1\}$. There are $2^n$ choices at each layer and $2^{nL}$ sequences in all.
\end{definition}

The word formal means that $D$ is chosen on paper; no input is claimed to produce it. For a fixed $D$ define the formal preactivation matrices by
\begin{equation}\label{eq:HD}
H^D_1=W_1,\qquad H^D_\ell=W_\ell D_{\ell-1}H^D_{\ell-1}\quad(2\le\ell\le L),\qquad\text{so}\qquad H^D_\ell=W_\ell D_{\ell-1}W_{\ell-1}\cdots D_1W_1,
\end{equation}
and the formal Jacobian
\begin{equation}\label{eq:JD}
J_D=D_LH^D_L=D_LW_LD_{L-1}W_{L-1}\cdots D_1W_1 .
\end{equation}
Compare with \cref{eq:jac}: the shape is identical, and the whole difference lies in where the masks came from.

To say ``the signs of a vector $y$ agree with the mask $D_\ell$'' with one inequality, put $E^D_\ell=2D_\ell-I_n$, a diagonal matrix with entries $+1$ where the mask is $1$ and $-1$ where it is $0$. Then
\[
E^D_\ell\,y>0\ \text{coordinatewise}\quad\Longleftrightarrow\quad y_i>0\ \text{where }(D_\ell)_{ii}=1\ \text{ and }\ y_i<0\ \text{where }(D_\ell)_{ii}=0 .
\]

\begin{definition}[strict region]
The strict region of a formal mask sequence $D$ is
\begin{equation}\label{eq:RD}
R_D=\bigl\{s\in\R^n:\ E^D_\ell H^D_\ell s>0\ \text{coordinatewise for every }1\le\ell\le L\bigr\}.
\end{equation}
\end{definition}

Every condition in \cref{eq:RD} is a strict homogeneous linear inequality in $s$, so $R_D$ is an open convex cone, possibly empty. It is open because strict inequalities survive small perturbations, convex because the inequalities are linear, and a cone because they are homogeneous.

\begin{lemma}[consistency]\label{lem:consistency}
If $s\in R_D$ then the forward pass at $s$ has mask $D_\ell$ at every layer, and
\[
h_\ell(s)=H^D_\ell s,\qquad x_\ell(s)=D_\ell H^D_\ell s,\qquad 1\le\ell\le L .
\]
In particular $F^{(\alpha)}_{n,L}(s)=J_Ds$ and $DF^{(\alpha)}_{n,L}(s)=J_D$. Conversely, an input at which every preactivation coordinate is nonzero lies in exactly one strict region.
\end{lemma}

\begin{proof}
Induction on the layer. At layer one, $h_1(s)=W_1s=H^D_1s$ by definition. Since $s\in R_D$, $E^D_1H^D_1s>0$ coordinatewise: where $(D_1)_{ii}=1$ this says $(H^D_1s)_i>0$ and ReLU keeps the coordinate; where $(D_1)_{ii}=0$ it says $(H^D_1s)_i<0$ and ReLU sends it to zero. Coordinate by coordinate, $x_1(s)=\phi(h_1(s))=D_1H^D_1s$. Suppose the claim through layer $\ell-1$. Then $h_\ell(s)=W_\ell x_{\ell-1}(s)=W_\ell D_{\ell-1}H^D_{\ell-1}s=H^D_\ell s$ by \cref{eq:HD}, and the inequality $E^D_\ell H^D_\ell s>0$ says the signs of $h_\ell(s)$ are those encoded by $D_\ell$, so $x_\ell(s)=D_\ell H^D_\ell s$. At the last layer $F(s)=x_L(s)=D_LH^D_Ls=J_Ds$. The strict inequalities hold in a neighbourhood of $s$, so the network is the linear map $J_D$ there and its derivative at $s$ is $J_D$.

For the converse, take an input with all preactivation coordinates nonzero and define $(D_\ell)_{ii}=\one_{\{h_{\ell i}(s)>0\}}$. The induction just made shows $h_\ell(s)=H^D_\ell s$ and every inequality of \cref{eq:RD} holds, so $s\in R_D$; and the sign of each nonzero $h_{\ell i}(s)$ determines $(D_\ell)_{ii}$, so $D$ is unique.
\end{proof}

\subsection{Realizability as a subspace meeting an orthant}

Stack the formal preactivations:
\begin{equation}\label{eq:AD}
A_D=\begin{pmatrix}H^D_1\\ H^D_2\\ \vdots\\ H^D_L\end{pmatrix}\in\R^{nL\times n},
\end{equation}
so that $A_Ds$ lists all $L$ formal preactivation vectors of $s$ on top of one another. Let
\[
O_D=\bigl\{(y_1,\dots,y_L)\in(\R^n)^L:\ E^D_\ell y_\ell>0\ \text{for every }\ell\bigr\},
\]
one open coordinate orthant of $\R^{nL}$. Then, directly from the definitions,
\begin{equation}\label{eq:realize}
R_D\ne\varnothing\quad\Longleftrightarrow\quad\im(A_D)\cap O_D\ne\varnothing .
\end{equation}
This is the geometric form of realizability that the symmetry argument of \cref{sec:gauge} works with: the mask $D$ is realized by some input exactly when the image of $A_D$, a subspace of dimension at most $n$ in a space of dimension $nL$, touches one specified orthant.

The event $\{R_D\ne\varnothing\}$ is measurable. Write $M=nL$, let $\tau_D\in\{\pm1\}^M$ be the sign vector of $O_D$, and for $A\in\R^{M\times n}$ put $\Psi_D(A)=\max_{\norm s=1}\min_i(\tau_D)_i(As)_i$, a maximum of a continuous function over the compact sphere. If $As\in O_D$ for some $s$ then $s\ne0$ and, by homogeneity, $\Psi_D(A)>0$; conversely a maximizing unit vector shows $\Psi_D(A)>0$ implies $\im A\cap O_D\ne\varnothing$. And $\abs{\Psi_D(A)-\Psi_D(B)}\le\sup_{\norm s=1}\norm{(A-B)s}_\infty$, so $\Psi_D$ is continuous and $\{\Psi_D>0\}$ is open. Since $A_D$ is a measurable function of the weights, so is the indicator of realizability.

\subsection{The Lipschitz constant of a piecewise-linear map}

\begin{lemma}\label{lem:pl}
Let $F:\R^n\to\R^m$ be continuous, and suppose a finite polyhedral complex covers $\R^n$ such that on every full-dimensional cell $C$ the map $F$ agrees with a linear map $M_C$. Then
\[
\Lip_2(F)=\max_C\op{M_C}.
\]
\end{lemma}

The Lipschitz constant compares points that may lie in different cells, while the operator norms are local to single cells. The proof shows that no stretching accumulates across cells, because the pieces of a straight segment add up to exactly its length.

\begin{proof}
Let $M=\max_C\op{M_C}$. Fix $s,t$ and parameterize the segment $\gamma(u)=s+u(t-s)$, $0\le u\le 1$. A finite complex has finitely many faces, each cut out by finitely many linear equations and inequalities. Thus the parameter interval has a finite subdivision on whose open subintervals the minimal face containing the segment is constant: $[0,1]$ splits as $0=u_0<u_1<\dots<u_r=1$ with the interior of each $\gamma([u_{j-1},u_j])$ inside one face. Choose a full-dimensional cell $C_j$ whose closure contains that face. On the interior of $C_j$, $F=M_{C_j}$; both sides are continuous, so they agree on the closure, in particular at the two endpoints of the subsegment. Therefore
\begin{align*}
\norm{F(t)-F(s)}&\le\sum_{j=1}^r\norm{F(\gamma(u_j))-F(\gamma(u_{j-1}))}
=\sum_{j=1}^r\norm{M_{C_j}\bigl(\gamma(u_j)-\gamma(u_{j-1})\bigr)}\\
&\le\sum_{j=1}^r\op{M_{C_j}}\norm{\gamma(u_j)-\gamma(u_{j-1})}
\le M\sum_{j=1}^r(u_j-u_{j-1})\norm{t-s}=M\norm{t-s}.
\end{align*}
The last equality is the point: along a straight line with constant velocity the displacements are $(u_j-u_{j-1})(t-s)$, and the parameter increments sum to one, so the pieces cannot accumulate. Hence $\Lip_2(F)\le M$.

For the reverse inequality choose a cell $C$ with $\op{M_C}=M$, an interior point $s\in C$, and a unit $v$ with $\norm{M_Cv}=M$. For small $h$, $s+hv\in C$, so $\norm{F(s+hv)-F(s)}/\abs h=\norm{M_Cv}=M$, and $\Lip_2(F)\ge M$.
\end{proof}

\subsection{Why the network has finitely many pieces, and the nondegeneracy event}

Each coordinate of $W_1s$ changes sign only across a hyperplane; on a region of fixed signs the first layer is multiplication by a fixed diagonal $0/1$ matrix, hence linear. Repeating the argument layer by layer refines $\R^n$ into finitely many polyhedral cells on which all signs are fixed and the network is linear; there are at most $2^{nL}$ sign sequences, so finitely many pieces.

One technical point: a preactivation functional could vanish identically on a cell, in which case the interior of that cell is not a strict region. With probability one this happens only after the formal trajectory has already died. Fix $D$, set $B^D_0=I_n$ and $B^D_{\ell-1}=D_{\ell-1}H^D_{\ell-1}$ for $\ell\ge 2$, so $H^D_\ell=W_\ell B^D_{\ell-1}$. Condition on $W_1,\dots,W_{\ell-1}$; then $B^D_{\ell-1}$ is fixed and row $i$ of $H^D_\ell$ is $w_i^TB^D_{\ell-1}$ with $w_i$ the $i$-th row of $W_\ell$. If $B^D_{\ell-1}\ne 0$, the map $w\mapsto w^TB^D_{\ell-1}$ is not zero, its kernel is a proper subspace, and a Gaussian vector lies in a proper subspace with probability zero. There are finitely many $(D,\ell,i)$, so almost surely: whenever $B^D_{\ell-1}\ne 0$, every row of $H^D_\ell$ is a nonzero functional. If instead $B^D_{\ell-1}=0$ then $H^D_\ell=0$ and everything after it, including $J_D$, is zero: a dead formal piece.

\begin{lemma}[exact reduction]\label{lem:reduction}
With probability one over the weights,
\begin{equation}\label{eq:reduction}
K_{n,L}(\alpha)=\max\bigl(\{0\}\cup\{\op{J_D}:R_D\ne\varnothing\}\bigr).
\end{equation}
Consequently, for every $t\ge 0$,
\begin{equation}\label{eq:reduction-event}
\{K_{n,L}(\alpha)>t\}\subseteq\bigcup_D\{R_D\ne\varnothing,\ \op{J_D}>t\}.
\end{equation}
\end{lemma}

\begin{proof}
Work on the nondegeneracy event. Refine input space into the finite complex on whose full-dimensional cells the network is linear. Take one such cell. If its formal forward map has not died, every preactivation functional used to define the next sign split is nonzero, its zero set is a hyperplane with empty interior, and the interior of each resulting full-dimensional subcell has strict signs at every layer; by \cref{lem:consistency} that subcell lies in a nonempty $R_D$ and the network's linear map there is $J_D$. If the formal map dies, bias-freeness makes every later output zero on the cell, so its linear map is the zero map. Thus every full-dimensional linear piece has matrix $J_D$ for a nonempty strict region or $0$, and \cref{lem:pl} gives \cref{eq:reduction}. If $K_{n,L}(\alpha)>t\ge 0$, the maximum exceeds $t$, the zero cannot be responsible, so some $D$ has $R_D\ne\varnothing$ and $\op{J_D}>t$.
\end{proof}

The lemma turns a supremum over all pairs of inputs into a maximum over finitely many objects, each attached to a formal mask. Boundary points between cells do not matter: the segment proof of \cref{lem:pl} is the reason.

\section{The symmetry, the realizable masks, and the counting}
\label{sec:gauge}

After \cref{lem:reduction} the problem is a maximum over the $2^{nL}$ formal masks, each contributing only if it is realized. Two things are still entangled: whether $D$ is realized, and how large $\op{J_D}$ is, are both functions of the same weights. This section separates them by an exact identity, not by an independence assumption.

\subsection{The gauge group}

For each layer let $\Sigma_\ell$ be a diagonal matrix with diagonal entries in $\{-1,+1\}$, and set $\Sigma_0=I_n$. Define transformed weights
\begin{equation}\label{eq:gauge}
W^\Sigma_1=\Sigma_1W_1,\qquad W^\Sigma_\ell=\Sigma_\ell W_\ell\Sigma_{\ell-1}\quad(2\le\ell\le L).
\end{equation}
Because $\Sigma_{\ell-1}^{-1}=\Sigma_{\ell-1}$, the same formula reads as a change of sign of the hidden coordinates at layer $\ell-1$: the sign inserted on the left of one layer is cancelled on the right of the next. The tuples $\Sigma=(\Sigma_1,\dots,\Sigma_L)$ form a group $\gauge$ with $\abs\gauge=2^{nL}$ elements. Write $S_\Sigma=\diag(\Sigma_1,\dots,\Sigma_L)\in\R^{nL\times nL}$.

\begin{example}
With $L=2$, $W^\Sigma_1=\Sigma_1W_1$ and $W^\Sigma_2=\Sigma_2W_2\Sigma_1$. The factor $\Sigma_1$ flips the signs of the first-layer preactivations; the matching factor on the right of $W_2$ cancels it inside the product; only $\Sigma_2$ survives in the formal Jacobian.
\end{example}

\begin{lemma}[gauge covariance]\label{lem:covariance}
Fix a formal mask sequence $D$. For every $\Sigma\in\gauge$,
\begin{equation}\label{eq:cov}
H^D_\ell(W^\Sigma)=\Sigma_\ell H^D_\ell(W),\qquad A_D(W^\Sigma)=S_\Sigma A_D(W),\qquad J_D(W^\Sigma)=\Sigma_LJ_D(W).
\end{equation}
Moreover $(W^\Sigma_1,\dots,W^\Sigma_L)\eqd(W_1,\dots,W_L)$.
\end{lemma}

\begin{proof}
At layer one, $H^D_1(W^\Sigma)=\Sigma_1W_1=\Sigma_1H^D_1(W)$. If the first identity holds at layer $\ell-1$, then
\begin{align*}
H^D_\ell(W^\Sigma)&=W^\Sigma_\ell D_{\ell-1}H^D_{\ell-1}(W^\Sigma)=\Sigma_\ell W_\ell\Sigma_{\ell-1}D_{\ell-1}\Sigma_{\ell-1}H^D_{\ell-1}(W)\\
&=\Sigma_\ell W_\ell D_{\ell-1}H^D_{\ell-1}(W)=\Sigma_\ell H^D_\ell(W),
\end{align*}
because diagonal matrices commute and $\Sigma_{\ell-1}^2=I$. Stacking gives the second identity; the third follows since $D_L$ and $\Sigma_L$ commute. For the law: each entry of $W^\Sigma_\ell$ is the corresponding entry of $W_\ell$ times a fixed sign, a centred Gaussian has the same law after a sign change, and the transformation is a bijective coordinatewise sign change, so independence within and between matrices is preserved.
\end{proof}

A nonnegative measurable function $F(W_1,\dots,W_L)$ is gauge invariant if $F(W^\Sigma)=F(W)$ for all $\Sigma$. Since $\Sigma_L$ is orthogonal, $\op{J_D(W^\Sigma)}=\op{\Sigma_LJ_D(W)}=\op{J_D(W)}$, and likewise $\norm{J_D(W^\Sigma)v}=\norm{J_D(W)v}$ for a fixed vector $v$: every nonnegative function of $\op{J_D}$ or of $\norm{J_Dv}$ is gauge invariant. The gauge changes the coordinates in which hidden representations are written; the internal signs cancel from the matrix product, and at the output only an orthogonal sign matrix remains, which cannot change a Euclidean norm.

\subsection{How many orthants can a subspace meet}

An open coordinate orthant of $\R^M$ is $O_\tau=\{x:\tau_ix_i>0\ \forall i\}$ for a sign vector $\tau\in\{\pm1\}^M$; there are $2^M$ of them. For a linear subspace $U\subseteq\R^M$ let
\[
r(U)=\#\{\text{open orthants that meet }U\}.
\]
If $U$ lies in a coordinate hyperplane then $r(U)=0$. The fact that drives the proof is that a $d$-dimensional subspace cannot pass through all $2^M$ orthants when $d\ll M$; it has only $d$ degrees of freedom and cannot choose all $M$ signs independently.

Assume no coordinate functional vanishes identically on $U$. The $M$ hyperplanes $\{x_i=0\}$ restrict to central hyperplanes in $U$ and cut $U$ into connected open cones. Two points of the same component have the same signs, since a sign change requires passing through zero; conversely the points of $U$ with a fixed strict sign pattern form the convex, hence connected, set $U\cap O_\tau$. So $r(U)$ is the number of regions into which the coordinate hyperplanes cut $U$, and the question becomes the classical count of regions of a central arrangement.

Let $R(M,d)$ be the largest number of regions that $M$ central hyperplanes can form in a space of dimension at most $d$. Adding the $M$-th hyperplane to an arrangement of $M-1$: a region it does not touch stays one region; a region it cuts becomes two, and the cut regions correspond one to one to the regions of the arrangement induced on the new hyperplane, which has dimension at most $d-1$ and carries at most $M-1$ central hyperplanes. Therefore
\begin{equation}\label{eq:recur}
R(M,d)\le R(M-1,d)+R(M-1,d-1),\qquad R(M,1)\le 2,\qquad R(1,d)\le 2 .
\end{equation}
Define $\widetilde R(M,d)=2\sum_{j=0}^{d-1}\binom{M-1}{j}$, binomials outside their range being zero. Pascal's identity $\binom{M-1}{j}=\binom{M-2}{j}+\binom{M-2}{j-1}$ gives
\[
\widetilde R(M-1,d)+\widetilde R(M-1,d-1)=2\sum_{j=0}^{d-1}\Bigl[\binom{M-2}{j}+\binom{M-2}{j-1}\Bigr]=2\sum_{j=0}^{d-1}\binom{M-1}{j}=\widetilde R(M,d),
\]
and $\widetilde R(M,1)=\widetilde R(1,d)=2$, so induction on \cref{eq:recur} gives $R(M,d)\le\widetilde R(M,d)$.

\begin{lemma}[orthant count]\label{lem:cover}
Let $U\subseteq\R^M$ be a linear subspace of dimension $d\le n$. Then $r(U)\le 2\sum_{j=0}^{d-1}\binom{M-1}{j}\le 2\sum_{j=0}^{n-1}\binom{M-1}{j}$.
\end{lemma}

\begin{proof}
If $U$ is inside a coordinate hyperplane, $r(U)=0$. Otherwise the region count above gives the first inequality, and $d\le n$ gives the second.
\end{proof}

For the network $M=nL$ and $U=\im(A_D)$ has dimension at most $n$. Define
\begin{equation}\label{eq:CnL}
C_{n,L}=2\sum_{j=0}^{n-1}\binom{nL-1}{j},\qquad q_{n,L}=\frac{C_{n,L}}{2^{nL}},
\end{equation}
so that $r(\im A_D)\le C_{n,L}$ always. Two values are exact: $C_{n,1}=2^n$ (every orthant of $\R^n$ is available to a full-rank $W_1$) and $C_{n,2}=2^{2n-1}$ (half of $2^{2n-1}$ binomials of an odd row).

\begin{lemma}[entropy bound]\label{lem:C-entropy}
For $L\ge 2$, $C_{n,L}\le 2\exp\bigl(nL\,H(1/L)\bigr)$.
\end{lemma}

\begin{proof}
Apply \cref{lem:entropy} with $N=nL-1$ and $k=n-1$: since $k/N\le 1/L\le 1/2$ and $H$ is increasing on $[0,1/2]$, $\sum_{j\le n-1}\binom{nL-1}{j}\le\exp((nL-1)H(\tfrac{n-1}{nL-1}))\le\exp(nLH(1/L))$.
\end{proof}

The ambient space has $2^{nL}=e^{nL\log 2}$ orthants; the subspace reaches at most $e^{nLH(1/L)}$ of them, and $LH(1/L)=\log L+1+O(1/L)$ grows like a logarithm of the depth, not linearly. This is the whole saving.

\begin{remark}[when the count is attained]
If the $M$ central hyperplanes are in general position in a $d$-dimensional space (every $k\le d$ of the defining functionals independent), the recurrence \cref{eq:recur} is an equality and the region count is exactly $2\sum_{j<d}\binom{M-1}{j}$: removing the last hyperplane leaves an arrangement in general position whose restriction to the removed hyperplane is again in general position, so every region it cuts is split in two and the count is additive. For the all-active mask this happens almost surely (\cref{sec:sharp}), which is how the exact realizability probability of that mask is computed and checked in \cref{sec:numerics}.
\end{remark}

\subsection{The realizable masks: the exact orbit identity}

\begin{proposition}[weighted gauge--orthant identity]\label{prop:orbit}
Fix a formal mask sequence $D$ and let $F=F(W_1,\dots,W_L)\ge 0$ be measurable and gauge invariant. Then
\begin{equation}\label{eq:orbit}
\E\bigl[F\,\one_{\{R_D\ne\varnothing\}}\bigr]=\E\Bigl[F\,\frac{r(\im A_D)}{2^{nL}}\Bigr].
\end{equation}
Consequently
\begin{equation}\label{eq:orbit-bound}
\E\bigl[F\,\one_{\{R_D\ne\varnothing\}}\bigr]\le q_{n,L}\,\E F .
\end{equation}
\end{proposition}

\begin{proof}
By \cref{eq:realize}, $\one_{\{R_D\ne\varnothing\}}=\one_{\{\im A_D(W)\cap O_D\ne\varnothing\}}$. Fix one $\Sigma\in\gauge$. The transformed weights have the same joint law as $W$, so
\[
\E\bigl[F(W)\one_{\{\im A_D(W)\cap O_D\ne\varnothing\}}\bigr]=\E\bigl[F(W^\Sigma)\one_{\{\im A_D(W^\Sigma)\cap O_D\ne\varnothing\}}\bigr].
\]
Gauge invariance gives $F(W^\Sigma)=F(W)$, and gauge covariance gives the identity $\im A_D(W^\Sigma)=S_\Sigma\im A_D(W)$. Hence
\[
\E\bigl[F\one_{\{R_D\ne\varnothing\}}\bigr]=\E\bigl[F(W)\one_{\{S_\Sigma\im A_D(W)\cap O_D\ne\varnothing\}}\bigr]\qquad\text{for every }\Sigma\in\gauge .
\]
Average over the $2^{nL}$ elements of the group:
\begin{equation}\label{eq:avg}
\E\bigl[F\one_{\{R_D\ne\varnothing\}}\bigr]=\E\Bigl[F(W)\,\frac{1}{2^{nL}}\sum_{\Sigma\in\gauge}\one_{\{S_\Sigma U\cap O_D\ne\varnothing\}}\Bigr],\qquad U=\im A_D(W).
\end{equation}
Now fix the realized $U$. Since $S_\Sigma^{-1}=S_\Sigma$, $S_\Sigma U\cap O_D\ne\varnothing$ if and only if $U\cap S_\Sigma O_D\ne\varnothing$. As $\Sigma$ ranges over the group, $S_\Sigma O_D$ ranges over all $2^{nL}$ open orthants exactly once: coordinate by coordinate, the fixed sign of $O_D$ is multiplied by an independently chosen sign, and every sign vector arises from exactly one $\Sigma$. Therefore
\[
\sum_{\Sigma\in\gauge}\one_{\{S_\Sigma U\cap O_D\ne\varnothing\}}=\#\{\text{orthants meeting }U\}=r(U),
\]
and \cref{eq:avg} is \cref{eq:orbit}. Finally $r(\im A_D)\le C_{n,L}$ by \cref{lem:cover}, and $F\ge 0$, so the right side of \cref{eq:orbit} is at most $(C_{n,L}/2^{nL})\E F$.
\end{proof}

Read the identity once more. On the left is the realizability of one formal mask, weighted by any gauge-invariant functional. On the right the realizability event has disappeared; in its place stands a deterministic-looking ratio, the number of orthants the random subspace $\im A_D$ meets divided by the number available. The identity is exact. It does not say that realizability and $F$ are independent: the random count $r(\im A_D)$ may well be correlated with $F$. What produces the useful inequality is that the count is bounded by a deterministic number, \cref{lem:cover}. The average over the gauge orbit is the only place the sign symmetry is used, and it is used exactly, not approximately.

This is also why the argument is unconditional: no input is chosen, no activation pattern is conditioned on, no witness input is constructed, and the law of the weights is never tilted. A mask is fixed on paper, a symmetry is averaged, and then all masks are summed.

\subsection{The bridge to an independent uniform mask}

Let $D$ now be random, independent of the weights, and uniform on the $2^{nL}$ formal mask sequences.

\begin{corollary}[unconditional bridge]\label{cor:bridge}
For every $t\ge 0$ and every $r>0$,
\begin{align}
\PP\bigl(K_{n,L}(\alpha)>t\bigr)&\le C_{n,L}\,\PP\bigl(\op{J_D}>t\bigr),\label{eq:bridge-tail}\\
\E\bigl[K_{n,L}(\alpha)^r\bigr]&\le C_{n,L}\,\E\op{J_D}^r .\label{eq:bridge-moment}
\end{align}
\end{corollary}

\begin{proof}
By \cref{lem:reduction} and the union bound, $\PP(K>t)\le\sum_D\PP(R_D\ne\varnothing,\op{J_D}>t)$. For fixed $D$ apply \cref{prop:orbit} with $F=\one_{\{\op{J_D}>t\}}$, which is nonnegative and gauge invariant: $\PP(R_D\ne\varnothing,\op{J_D}>t)\le q_{n,L}\PP(\op{J_D}>t)$. Summing over the $2^{nL}$ fixed masks, $\sum_Dq_{n,L}\PP(\op{J_D}>t)=C_{n,L}\,2^{-nL}\sum_D\PP(\op{J_D}>t)=C_{n,L}\PP(\op{J_D}>t)$ with $D$ uniform and independent. For the moment, \cref{eq:reduction} gives pointwise $K^r\le\sum_D\one_{\{R_D\ne\varnothing\}}\op{J_D}^r$; take expectations and apply \cref{prop:orbit} with $F=\op{J_D}^r$.
\end{proof}

The left side of \cref{eq:bridge-tail} is the full adaptive supremum over every input. The right side is the Jacobian of a mask chosen by a coin toss, independent of the weights. The entire cost of passing from one to the other is the factor $C_{n,L}$, of exponential size $e^{nLH(1/L)}$ against the $2^{nL}$ masks that a naive union bound would have paid for.

\subsection{Testing the identity}

An exact identity can be tested, and \cref{eq:orbit} was, at width two, where both sides can be computed without sampling error in the direction variable: in the plane the sign pattern of $A_Ds$ is constant on each arc between the angles at which a row of $A_D$ vanishes, so enumerating the arcs gives the exact set of orthants met by $\im A_D$, hence $r(\im A_D)$ exactly and the realizability indicator exactly. \Cref{tab:orbit} reports the paired Monte Carlo test over the weights for $F\equiv 1$, which turns \cref{eq:orbit} into
\[
\PP(R_D\ne\varnothing)=\E\,r(\im A_D)/2^{nL},
\]
and for the bounded functional $F=\min(\op{J_D},1)$. The initial calculations illustrate two limitations of numerical checks. Sweeping eight thousand directions round the circle instead of enumerating the arcs missed thin arcs and pulled both sides down by about six percent, consistently; sampling a set of arcs is not enumerating it, and the error appeared as a bias rather than noise. Testing with $F=\op{J_D}^2$ produced a paired difference two and a half estimated standard errors from zero at depth three. That observation alone establishes neither a failure of the identity nor undercoverage by the standard-error estimate. An unbounded functional can be sensitive to rare large draws, so a bounded functional provides a useful additional check. The identity holds for every nonnegative gauge-invariant $F$; with the bounded choices used here, every ratio is within two percent of one and every paired $z$-score within $1.3$ of zero.

\section{The machine: nets, exact moments, and one scalar variable}
\label{sec:moments}

\Cref{cor:bridge} leaves a Jacobian $J_D$ with a mask independent of the weights. This section bounds its operator norm by a finite net, computes the moments of $\norm{J_Dv}$ for a fixed $v$ exactly, and reduces everything to one scalar random variable.

\subsection{A net of the sphere}

\begin{definition}
A finite set $\net\subseteq S^{n-1}$ is an $\varepsilon$-net if for every $x\in S^{n-1}$ there is $v\in\net$ with $\norm{x-v}\le\varepsilon$.
\end{definition}

\begin{lemma}\label{lem:net}
There is a $\tfrac12$-net $\net\subseteq S^{n-1}$ with $\abs\net\le 5^n$, and for every linear $T:\R^n\to\R^m$,
\begin{equation}\label{eq:net}
\op T\le 2\max_{v\in\net}\norm{Tv}.
\end{equation}
\end{lemma}

\begin{proof}
Take a maximal $\tfrac12$-separated subset $\net$ of the sphere: distinct $u,v\in\net$ have $\norm{u-v}>\tfrac12$ and no point of the sphere can be added. Maximality gives the net property, for a point farther than $\tfrac12$ from all of $\net$ could be added. The balls of radius $\tfrac14$ about the points of $\net$ are disjoint, since a common point $z$ would give $\norm{u-v}\le\norm{u-z}+\norm{z-v}\le\tfrac12$; and each lies in the ball of radius $\tfrac54$ about the origin, since $\norm z\le\norm{z-v}+\norm v\le\tfrac14+1$. Comparing volumes, which scale as the $n$-th power of the radius, $\abs\net\,(\tfrac14)^n\le(\tfrac54)^n$, so $\abs\net\le 5^n$.

For \cref{eq:net}, let $x$ be a unit vector and $v\in\net$ with $\norm{x-v}\le\tfrac12$. Then $\norm{Tx}\le\norm{Tv}+\norm{T(x-v)}\le\max_{u\in\net}\norm{Tu}+\tfrac12\op T$. Take the supremum over $x$ and move the last term to the left.
\end{proof}

Raising \cref{eq:net} to the power $2p$ and using $\max_ja_j\le\sum_ja_j$ for nonnegative $a_j$,
\begin{equation}\label{eq:net-moment}
\op T^{2p}\le 4^p\sum_{v\in\net}\norm{Tv}^{2p}\qquad(p>0).
\end{equation}
Only a net of input directions is needed, because $\norm{Tv}$ is a vector norm; a bilinear estimate would need a second net, and the extra factor is unnecessary here. The net never looks for activation regions: it is applied to the fixed linear map $J_D$, so an arbitrarily narrow realizable cone is never missed.

\subsection{Exact moments of $\norm{J_Dv}$ for a fixed $v$}

Fix a formal mask sequence $D$ and let $S_\ell=\{i:(D_\ell)_{ii}=1\}$ with $m_\ell=\abs{S_\ell}$ active coordinates. Let $R_\ell:\R^n\to\R^{m_\ell}$ keep exactly the coordinates in $S_\ell$; its transpose inserts a vector into those coordinates with zeros elsewhere, so $D_\ell=R_\ell^TR_\ell$, $R_\ell R_\ell^T=I_{m_\ell}$, and $\norm{R_\ell^Ty}=\norm y$. With $R_0=I_n$, $m_0=n$, define the active blocks
\[
G_1=R_1W_1,\qquad G_\ell=R_\ell W_\ell R_{\ell-1}^T\quad(2\le\ell\le L),
\]
of size $m_\ell\times m_{\ell-1}$, obtained by selecting rows and columns of $W_\ell$; so all their entries are independent $\mathcal N(0,\alpha/n)$ and $G_1,\dots,G_L$ are independent. Inserting $D_\ell=R_\ell^TR_\ell$ into \cref{eq:JD},
\[
J_D=R_L^T(R_LW_LR_{L-1}^T)(R_{L-1}W_{L-1}R_{L-2}^T)\cdots(R_1W_1)=R_L^TG_LG_{L-1}\cdots G_1,
\]
and since $R_L^T$ preserves norms,
\begin{equation}\label{eq:JDv}
\norm{J_Dv}=\norm{G_LG_{L-1}\cdots G_1v}.
\end{equation}
If some $m_\ell=0$ the block has no rows and the product is zero; the formulas below remain valid with $\chi^2_0=0$.

One Gaussian block acting on one vector: let $G$ be $m\times k$ with independent $\mathcal N(0,\alpha/n)$ entries and $y\in\R^k$ fixed. Row $i$ of $Gy$ is $\sum_jG_{ij}y_j$, centred Gaussian with variance $\frac\alpha n\norm y^2$; different rows use disjoint entries, so the coordinates are independent. Hence $Gy\eqd\sqrt{\alpha/n}\,\norm y\,Z$ with $Z\sim\mathcal N(0,I_m)$, and
\begin{equation}\label{eq:one-block}
\norm{Gy}^2\eqd\frac\alpha n\norm y^2\chi^2_m .
\end{equation}

\begin{lemma}[exact radial factors]\label{lem:radial}
For every fixed unit vector $v\in\R^n$ and every $p>0$,
\begin{equation}\label{eq:radial}
\E\norm{J_Dv}^{2p}=\prod_{\ell=1}^L\E\Bigl(\frac\alpha n\chi^2_{m_\ell}\Bigr)^p .
\end{equation}
\end{lemma}

\begin{proof}
Set $v_0=v$ and $v_\ell=G_\ell v_{\ell-1}$, so $v_L=G_L\cdots G_1v$ and $\norm{J_Dv}=\norm{v_L}$ by \cref{eq:JDv}. Conditional on $v_{\ell-1}$, \cref{eq:one-block} gives $\norm{v_\ell}^2\eqd\norm{v_{\ell-1}}^2\frac\alpha n\chi^2_{m_\ell}$, so $\E[\norm{v_\ell}^{2p}\mid v_{\ell-1}]=\norm{v_{\ell-1}}^{2p}\E(\frac\alpha n\chi^2_{m_\ell})^p$, the factor being deterministic once $m_\ell$ is fixed. Take expectations (the tower property) and iterate from $\ell=L$ down to $1$; $\norm{v_0}=1$.
\end{proof}

The vectors $v_1,\dots,v_L$ are not independent, and no such claim is made. The product formula rests only on the conditional identity and on the independence of the new block $G_\ell$ from everything before it. The expectation depends on $D$ only through the active counts $m_1,\dots,m_L$ and not on the direction $v$.

\subsection{The scalar mixture}

Let $B_1,\dots,B_n$ be independent Bernoulli$(\tfrac12)$ and $Z_1,\dots,Z_n$ independent standard Gaussians, independent of the $B_i$, and put
\begin{equation}\label{eq:X}
X_n^{(\alpha)}=\frac\alpha n\sum_{i=1}^nB_iZ_i^2 .
\end{equation}
Conditional on $M_n=\sum_iB_i=m$, exactly $m$ Gaussian squares appear, so $\sum_iB_iZ_i^2\mid\{M_n=m\}\eqd\chi^2_m$, and $M_n\sim\mathrm{Bin}(n,\tfrac12)$. Therefore, for every $p>0$,
\begin{equation}\label{eq:mixture}
\E\bigl[(X_n^{(\alpha)})^p\bigr]=2^{-n}\sum_{m=0}^n\binom nm\E\Bigl(\frac\alpha n\chi^2_m\Bigr)^p .
\end{equation}
This is exactly the sum that appears when the fixed-mask formula \cref{eq:radial} is averaged over a uniform mask.

\subsection{The master moment inequality}

\begin{proposition}[master inequality]\label{prop:master}
For every $n,L\ge 1$, $\alpha>0$ and $p>0$,
\begin{equation}\label{eq:master}
\E\bigl[K_{n,L}(\alpha)^{2p}\bigr]\le C_{n,L}\,5^n\,4^p\,\Bigl(\E\bigl[(X_n^{(\alpha)})^p\bigr]\Bigr)^L,
\end{equation}
and consequently
\begin{equation}\label{eq:master-tail}
\PP\bigl(K_{n,L}(\alpha)>1\bigr)\le C_{n,L}\,5^n\,4^p\,\Bigl(\E\bigl[(X_n^{(\alpha)})^p\bigr]\Bigr)^L .
\end{equation}
\end{proposition}

\begin{proof}
Fix the net $\net$ of \cref{lem:net}. By \cref{eq:reduction}, pointwise $K^{2p}\le\sum_D\one_{\{R_D\ne\varnothing\}}\op{J_D}^{2p}$, and by \cref{eq:net-moment} applied to each $J_D$,
\[
K_{n,L}(\alpha)^{2p}\le 4^p\sum_D\sum_{v\in\net}\one_{\{R_D\ne\varnothing\}}\norm{J_Dv}^{2p}.
\]
For fixed $D$ and $v$ the functional $\norm{J_D(W)v}^{2p}$ is nonnegative and gauge invariant, so \cref{prop:orbit} gives $\E[\one_{\{R_D\ne\varnothing\}}\norm{J_Dv}^{2p}]\le(C_{n,L}/2^{nL})\E\norm{J_Dv}^{2p}$. Hence
\[
\E K_{n,L}(\alpha)^{2p}\le 4^p\frac{C_{n,L}}{2^{nL}}\sum_D\sum_{v\in\net}\E\norm{J_Dv}^{2p}\le 4^p5^n\frac{C_{n,L}}{2^{nL}}\sum_D\E\norm{J_Dv}^{2p},
\]
using that the expectation does not depend on $v$ (\cref{lem:radial}) and $\abs\net\le 5^n$. Group the masks by their active counts: there are $\binom nm$ diagonal masks with $m$ ones at one layer, so by \cref{eq:radial}
\[
\sum_D\E\norm{J_Dv}^{2p}=\prod_{\ell=1}^L\Bigl[\sum_{m=0}^n\binom nm\E\Bigl(\frac\alpha n\chi^2_m\Bigr)^p\Bigr]=\Bigl[2^n\,\E(X_n^{(\alpha)})^p\Bigr]^L
\]
by \cref{eq:mixture}. The factor $2^{nL}$ cancels exactly, which proves \cref{eq:master}. On the event $\{K>1\}$ one has $K^{2p}>1$, so Markov's inequality gives \cref{eq:master-tail}.
\end{proof}

The four factors and their sources:
\begin{center}
\begin{tabular}{ll}
\toprule
factor & source\\
\midrule
$C_{n,L}$ & the orthants an at-most-$n$-dimensional stacked subspace can meet\\
$5^n$ & the size of a $\tfrac12$-net of the input sphere\\
$4^p$ & $\op T\le 2\max_{v\in\net}\norm{Tv}$, raised to the power $2p$\\
$(\E X_n^p)^L$ & one exact radial moment for each of the $L$ independent layers\\
\bottomrule
\end{tabular}
\end{center}
The proof will choose $p$ proportional to $n$. Then every factor has an exponential rate in $n$, and the scalar moment's rate, which is negative when $\alpha<2$, is repeated $L$ times.

\subsection{The finite-width scalar moment bound}

\begin{lemma}\label{lem:mgf}
Let $Y=BZ^2$ with $B\sim\mathrm{Bernoulli}(\tfrac12)$ and $Z\sim\mathcal N(0,1)$ independent, and $S_n=\sum_{i=1}^nB_iZ_i^2$ a sum of $n$ independent copies. For $t<\tfrac12$,
\begin{equation}\label{eq:mgfS}
\E e^{tY}=\frac{1+(1-2t)^{-1/2}}{2},\qquad\E e^{tS_n}=\Bigl[\frac{1+(1-2t)^{-1/2}}{2}\Bigr]^n .
\end{equation}
Also $\E Y=\tfrac12$, $\E Y^2=\tfrac32$, $\Var(Y)=\tfrac54$.
\end{lemma}

\begin{proof}
$\E e^{tY}=\tfrac12e^0+\tfrac12\E e^{tZ^2}$ and \cref{eq:mgf}; independence gives the product. $\E Y=\E B\,\E Z^2=\tfrac12$, $\E Y^2=\E B\,\E Z^4=\tfrac32$.
\end{proof}

\begin{lemma}[finite-width scalar moment bound]\label{lem:h}
For every $n\ge 1$ and $0<\theta<\tfrac14$,
\begin{equation}\label{eq:hbound}
\E\bigl[(X_n^{(\alpha)})^{\theta n}\bigr]\le e^{n\,h_\alpha(\theta)},\qquad h_\alpha(\theta)=\theta\log\frac\alpha{2e}+\log\Bigl(\frac{1+(1-4\theta)^{-1/2}}2\Bigr).
\end{equation}
Moreover $h_\alpha(0)=0$, $h_\alpha'(0+)=\log(\alpha/2)$, and as $\theta\downarrow0$
\begin{equation}\label{eq:h-expansion}
h_\alpha(\theta)=\theta\log\frac\alpha2+\frac52\theta^2+O(\theta^3),
\end{equation}
with a remainder that does not depend on $\alpha$.
\end{lemma}

\begin{proof}
Put $p=\theta n$ and $t=2\theta<\tfrac12$. Since $X_n^{(\alpha)}=\frac\alpha nS_n$, \cref{lem:xp} gives
\[
\E(X_n^{(\alpha)})^p=\Bigl(\frac\alpha n\Bigr)^p\E S_n^p\le\Bigl(\frac\alpha n\Bigr)^p\Bigl(\frac p{et}\Bigr)^p\E e^{tS_n}=\Bigl(\frac\alpha n\cdot\frac n{2e}\Bigr)^{\theta n}\Bigl[\frac{1+(1-4\theta)^{-1/2}}2\Bigr]^n=e^{nh_\alpha(\theta)},
\]
using $p/(et)=\theta n/(2e\theta)=n/(2e)$ and \cref{eq:mgfS}. For the expansion, $(1-x)^{-1/2}=1+\tfrac12x+\tfrac38x^2+\tfrac5{16}x^3+O(x^4)$ with $x=4\theta$ gives $(1-4\theta)^{-1/2}=1+2\theta+6\theta^2+20\theta^3+O(\theta^4)$, so $\frac{1+(1-4\theta)^{-1/2}}2=1+\theta+3\theta^2+10\theta^3+O(\theta^4)$, and $\log(1+y)=y-\tfrac12y^2+\tfrac13y^3+O(y^4)$ with $y=\theta+3\theta^2+10\theta^3+O(\theta^4)$ gives $\log(\cdots)=\theta+\tfrac52\theta^2+O(\theta^3)$. Since $\theta\log\frac\alpha{2e}=\theta\log\frac\alpha2-\theta$, the $\pm\theta$ cancel and \cref{eq:h-expansion} follows. All dependence on $\alpha$ sits in the exact linear term.
\end{proof}

The derivative at zero decides the qualitative threshold: if $\alpha<2$ then $h_\alpha'(0+)<0$ and $h_\alpha(\theta)<0$ for all small positive $\theta$. At $\alpha=2$ the linear term vanishes and $h_2(\theta)=\tfrac52\theta^2+O(\theta^3)$; the coefficient $\tfrac52$ is the normalized variance $\Var(Y)/(2(\E Y)^2)=\tfrac{5/4}{2\cdot 1/4}$ of one masked Gaussian square, and it is where the constant $\sqrt{10}$ of the rate comes from. \Cref{sec:sharp} shows that the large-$n$ exponent of the moment is exactly this function to second order, so neither number is an artifact of \cref{lem:xp}.

\section{The theorems}
\label{sec:theorems}

\subsection{The finite-width theorem}

\begin{proof}[Proof of \cref{thm:finite}]
Take $p=\theta n$ in \cref{prop:master}. The first inequality in \cref{eq:main-finite} is Markov's inequality on the event $\{K>1\}$, as in the proof of \cref{eq:master-tail}. For the second, \cref{lem:h} gives $\E(X_n^{(\alpha)})^{\theta n}\le e^{nh_\alpha(\theta)}$, so
\[
\E K_{n,L}(\alpha)^{2\theta n}\le C_{n,L}\,5^n\,4^{\theta n}\,e^{nLh_\alpha(\theta)},
\]
and \cref{lem:C-entropy} bounds $C_{n,L}\le 2e^{nLH(1/L)}$ for $L\ge2$. Multiplying out,
\[
C_{n,L}5^n4^{\theta n}e^{nLh_\alpha(\theta)}\le 2\exp\bigl(n[LH(1/L)+\log5+\theta\log4+Lh_\alpha(\theta)]\bigr)=2e^{nB_L(\alpha,\theta)}.
\]
For $L=1$, $C_{n,1}=2^n=2\cdot2^{n-1}\le2e^{n\log2}$ and the same computation applies with $\log2$ in place of $LH(1/L)$.
\end{proof}

\Cref{cor:width} is the theorem with the best $\theta$. The function $\theta\mapsto-B_L(\alpha,\theta)$ is continuous on $(0,1/4)$, tends to $-\infty$ as $\theta\uparrow1/4$, where $(1-4\theta)^{-1/2}$ blows up, and tends to the negative number $-LH(1/L)-\log5$ as $\theta\downarrow0$; so whenever its supremum $c_L(\alpha)$ is positive it is attained at an interior $\theta$, and the corollary applies with that $\theta$.

\subsection{The exponent budget}
\label{sec:budget}

Written as a rate per unit width, the bound reads
\[
\frac1n\log\PP\bigl(K_{n,L}(\alpha)>1\bigr)\le\frac{\log2}n+
\underbrace{LH(1/L)}_{\text{realizable patterns}}+\underbrace{\log5}_{\text{sphere net}}+\underbrace{\theta\log4}_{\text{net to norm}}+\underbrace{Lh_\alpha(\theta)}_{\text{radial moment, $L$ layers}} .
\]
Three of the four terms do not move once $L$ is fixed, and the first grows only like $\log L$: by \cref{lem:LH} below, $LH(1/L)=\log L+1+O(1/L)$. The last term is $L$ times a number that is negative whenever $\alpha<2$ and $\theta$ is small. A negative multiple of $L$ eventually beats a logarithm of $L$; the depth at which it does is $L_0(\alpha)$. This is why the upper bound succeeds at large depth and only there: the realizability entropy costs order $\log L$ per unit width, against the $nL\log2$ that a union bound over all $2^{nL}$ masks would cost, and the radial contraction is order $L$.

\begin{lemma}\label{lem:LH}
$LH(1/L)=\log L-(L-1)\log(1-1/L)=\log L+1+O(1/L)$ as $L\to\infty$, and $LH(1/L)\le\log L+1$ for every $L\ge2$.
\end{lemma}

\begin{proof}
$LH(1/L)=-\log(1/L)-(L-1)\log(1-1/L)$. Since $-\log(1-x)=x+x^2/2+\dots\ge x$, $-(L-1)\log(1-1/L)\ge(L-1)/L$, and since $-\log(1-x)\le x/(1-x)$, $-(L-1)\log(1-1/L)\le(L-1)\cdot\frac{1/L}{1-1/L}=1$. So $\log L+1-1/L\le LH(1/L)\le\log L+1$.
\end{proof}

\subsection{Fixed depth, growing width; the depth threshold}

\begin{proof}[Proof of \cref{thm:main}(i)]
Fix $\alpha<2$. By \cref{eq:h-expansion}, $h_\alpha'(0+)=\log(\alpha/2)<0$, so there is $\theta\in(0,1/4)$, depending on $\alpha$ but not on $n$ or $L$, with $h_\alpha(\theta)=-\gamma<0$. By \cref{lem:LH} the budget satisfies
\[
B_L(\alpha,\theta)\le\log L+1+\log5+\theta\log4-L\gamma,
\]
which is negative for all $L\ge L_0(\alpha)$, a linear function eventually dominating a logarithm. For such $L$, \cref{thm:finite} gives $\PP(K_{n,L}(\alpha)>1)\le2e^{-c_{\alpha,L}n}$ with $c_{\alpha,L}=-B_L(\alpha,\theta)>0$, for every width $n$.
\end{proof}

The smallest depth at which the exact budget $\min_\theta B_L(\alpha,\theta)$ turns negative is a computable number; for $\alpha=1$ it is $237$, for $\alpha=1.9$ it is $52\,841$ (\cref{sec:numerics}). At depths below it the bound says nothing, and at depth one it cannot: $\alpha_1=1/4$.

\subsection{The other side}

\begin{lemma}[one fixed input]\label{lem:forward}
Fix $L$ and a deterministic $s\ne0$. Conditional on $x_{\ell-1}(s)\ne0$,
\begin{equation}\label{eq:ratio}
\frac{\norm{x_\ell(s)}^2}{\norm{x_{\ell-1}(s)}^2}\eqd\frac\alpha n\sum_{i=1}^nB_{\ell i}Z_{\ell i}^2=X_n^{(\alpha)},
\end{equation}
with the pairs $(B_{\ell i},Z_{\ell i})$ independent across $i$. Also, for every nonzero $s$, whether or not the network is differentiable there,
\begin{equation}\label{eq:K-lower}
K_{n,L}(\alpha)\ge\frac{\norm{F^{(\alpha)}_{n,L}(s)}}{\norm s}.
\end{equation}
\end{lemma}

\begin{proof}
Given the past through layer $\ell-1$, $x_{\ell-1}(s)$ is fixed and $W_\ell$ independent of it, so the coordinates of $W_\ell x_{\ell-1}(s)$ are independent centred Gaussians of variance $\frac\alpha n\norm{x_{\ell-1}(s)}^2$; ReLU keeps each with probability $\tfrac12$, independently of its magnitude (a centred Gaussian's sign and absolute value are independent), which is \cref{eq:ratio}. For \cref{eq:K-lower}, $F(0)=0$ by \cref{eq:homog}, so $K\ge\norm{F(s)-F(0)}/\norm{s-0}$.
\end{proof}

\begin{proof}[Proof of \cref{thm:above2}]
Fix $s=e_1$ and write $r_\ell=\norm{x_\ell}^2/\norm{x_{\ell-1}}^2$ when $x_{\ell-1}\ne0$. If no layer dies and every $r_\ell>1$, then $\norm{x_L}^2/\norm s^2=\prod_\ell r_\ell>1$ and $K_{n,L}(\alpha)>1$ by \cref{eq:K-lower}. So
\[
\{K_{n,L}(\alpha)\le1\}\subseteq\bigcup_{\ell=1}^L\{x_{\ell-1}\ne0,\ \tfrac\alpha nS^{(\ell)}_n\le1\}\cup\{\text{some layer dies}\},
\]
where $S^{(\ell)}_n=\sum_iB_{\ell i}Z_{\ell i}^2$ is the sum in \cref{eq:ratio}. A layer dies exactly when all $n$ of its Bernoulli variables are zero, in which case $S^{(\ell)}_n=0\le n/\alpha$; so the dying event is contained in the same union. By the union bound and \cref{eq:ratio}, $\PP(K\le1)\le L\,\PP(S_n\le n/\alpha)$ with $S_n$ as in \cref{lem:mgf}. Since $\E S_n=n/2$, $\Var S_n=5n/4$ and $n/\alpha<n/2$ for $\alpha>2$, Chebyshev's inequality gives
\[
\PP\Bigl(S_n\le\frac n\alpha\Bigr)\le\PP\Bigl(\abs{S_n-\tfrac n2}\ge\frac n2-\frac n\alpha\Bigr)\le\frac{5n/4}{n^2(\frac12-\frac1\alpha)^2}=\frac{5\alpha^2}{n(\alpha-2)^2}.
\]
Multiply by $L$.
\end{proof}

\Cref{thm:main}(ii) follows, since the bound tends to zero as $n\to\infty$ at fixed $L$ and $\alpha>2$. A Chernoff bound with the moment-generating function \cref{eq:mgfS} gives an exponentially small bound in place of $1/n$; the weaker form is enough for the theorem and has an explicit constant.

\begin{remark}
The reverse direction needs no largest singular value, no rare activation pattern, and no global search. The ordinary forward direction of one fixed input already expands, because each layer multiplies the squared norm by an average of $\alpha/2>1$ and the fluctuations are of order $n^{-1/2}$.
\end{remark}

\subsection{The depth limit}

\begin{proof}[Proof of \cref{thm:main}(iii)]
Part (ii) says that no $\alpha>2$ belongs to the set in \cref{eq:alphaL} for any $L$, so $\alpha_L\le2$ for every $L$. Fix $\alpha<2$; by part (i), for every $L\ge L_0(\alpha)$, $\PP(K_{n,L}(\alpha)\le1)\to1$, so $\alpha$ belongs to the set and $\alpha_L\ge\alpha$ for all such $L$. Hence $\liminf_L\alpha_L\ge\alpha$ for every $\alpha<2$, so $\liminf_L\alpha_L\ge2\ge\limsup_L\alpha_L$.
\end{proof}

Reading the statement carefully: for every $\alpha<2$ there is a depth $L_0(\alpha)$, inside the choice of $\alpha$, such that for every $L\ge L_0(\alpha)$ the probability statement holds as $n\to\infty$. The order of the quantifiers cannot be exchanged: no single depth works for every $\alpha<2$ at once, because $L_0(\alpha)\to\infty$ as $\alpha\uparrow2$.

\subsection{The quantitative rate}

Write $\alpha=2e^{-\delta}$. Then $\theta\log(\alpha/2)=-\delta\theta$, and by \cref{eq:h-expansion}
\begin{equation}\label{eq:h-delta}
h_{2e^{-\delta}}(\theta)=-\delta\theta+\frac52\theta^2+O(\theta^3),
\end{equation}
with a remainder uniform in $\delta$ near zero, since all dependence on $\alpha$ is in the exact linear term. Ignoring the remainder, $-\delta\theta+\frac52\theta^2$ is minimized at $\theta=\delta/5$ with value $-\delta^2/10$; so $L$ layers produce a negative exponent of about $-L\delta^2/10$, against a geometric cost of about $\log L$. Balancing, $L\delta^2/10\approx\log L$, or $\delta\approx\sqrt{10\log L/L}$. The proof makes this precise.

\begin{proof}[Proof of \cref{thm:main}(iv)]
Fix a constant $A>1+\log5$ and for large $L$ put
\[
\delta_L=\sqrt{\frac{10(\log L+A)}L},\qquad\alpha(L)=2e^{-\delta_L},\qquad\theta_L=\frac{\delta_L}5 .
\]
Since $\delta_L\to0$, $\theta_L\in(0,1/4)$ for all large $L$. Evaluate the budget of \cref{thm:finite} at $(\alpha(L),\theta_L)$, with $L$ fixed and $n$ free:
\begin{align*}
B_L(\alpha(L),\theta_L)&=LH(1/L)+\log5+\theta_L\log4+Lh_{\alpha(L)}(\theta_L)\\
&=\log L+1+O(1/L)+\log5+o(1)+L\Bigl[-\frac{\delta_L^2}5+\frac52\cdot\frac{\delta_L^2}{25}+O(\delta_L^3)\Bigr]\\
&=\log L+1+\log5-\frac{L\delta_L^2}{10}+O(L\delta_L^3)+o(1).
\end{align*}
By definition $L\delta_L^2/10=\log L+A$, while $L\delta_L^3=10^{3/2}(\log L+A)^{3/2}L^{-1/2}\to0$. Hence
\[
B_L(\alpha(L),\theta_L)=1+\log5-A+o(1),
\]
which is negative for all large $L$ because $A>1+\log5$. For such $L$, \cref{thm:finite} gives $\PP(K_{n,L}(\alpha(L))>1)\to0$ as $n\to\infty$, so $\alpha(L)$ belongs to the set defining $\alpha_L$ and
\[
\alpha_L\ge\alpha(L)=2\exp\Bigl[-\sqrt{\frac{10(\log L+A)}L}\Bigr]=2\exp\Bigl[-(\sqrt{10}+o(1))\sqrt{\frac{\log L}L}\Bigr],
\]
the last step because $\sqrt{1+A/\log L}=1+o(1)$.
\end{proof}

The rate is the balance between one number that comes from geometry, the entropy $\log L$ of the realizable patterns, and one that comes from probability, the curvature $\tfrac52$ of the moment exponent at the critical variance. The asymptotic form drops the $O(1/L)$, the $\theta_L\log4$, the $O(L\delta^3_L)$ and the choice of $A$, and is therefore only the large-$L$ shape of the bound; at a specific depth the number the proof delivers is the threshold of the exact budget,
\begin{equation}\label{eq:alpha-star}
\alpha^*_L=\sup\Bigl\{\alpha:\ \min_{0<\theta<1/4}B_L(\alpha,\theta)<0\Bigr\}\le\alpha_L,
\end{equation}
which \cref{sec:numerics} tabulates. At $L=100$ the exact budget gives $0.678$ while the asymptotic formula, read at $L=100$ with the $o(1)$ dropped, would say $1.01$, a number the proof does not deliver there.

\subsection{Exact scaling in the variance}

Let $G_1,\dots,G_L$ have independent $\mathcal N(0,1/n)$ entries and couple $W^{(\alpha)}_\ell=\sqrt\alpha\,G_\ell$. Each layer contributes one factor $\sqrt\alpha$, and by homogeneity of ReLU, $F^{(\alpha)}_{n,L}(s)=\alpha^{L/2}F^{(1)}_{n,L}(s)$, hence
\begin{equation}\label{eq:scaling}
K_{n,L}(\alpha)=\alpha^{L/2}K_{n,L}(1).
\end{equation}
The positive scalar changes no activation sign, so every mask is unchanged under the coupling. Stability is therefore monotone in $\alpha$: the set in \cref{eq:alphaL} is an interval $(0,\alpha_L)$ or $(0,\alpha_L]$, and the word threshold is literal. The identity also says that the whole question is the distribution of $K_{n,L}(1)$, and that $\PP(K_{n,L}(\alpha)\le1)=\PP(K_{n,L}(1)\le\alpha^{-L/2})$.

\subsection{Biases}
\label{sec:biases}

Add bias vectors to the network: $h_\ell(s)=W_\ell x_{\ell-1}(s)+b_\ell$, with $b_1,\dots,b_L\in\R^n$ random, independent of the weights, and each with independent coordinates whose laws are symmetric about zero (Gaussian biases of any variance are the standard case). Everything above goes through with the changes listed here, and the bound of \cref{thm:finite} holds word for word.

\begin{theorem}[finite width, with biases]\label{thm:bias}
For the biased network, for every $n\ge1$, $L\ge2$, $\alpha>0$ and $\theta\in(0,1/4)$, $\PP(K_{n,L}(\alpha)>1)\le2\exp(nB_L(\alpha,\theta))$ with the same budget $B_L$ as in \cref{thm:finite}; and for $\alpha>2$, $\PP(K_{n,L}(\alpha)\le1)\le5\alpha^2L/(n(\alpha-2)^2)$. Consequently \cref{thm:main} holds for the biased network.
\end{theorem}

\begin{proof}
\emph{Regions.} For a formal mask sequence $D$ the formal preactivations are affine: $h^D_\ell(s)=H^D_\ell s+c^D_\ell$ with the same matrices $H^D_\ell$ as in \cref{eq:HD} and offsets $c^D_1=b_1$, $c^D_\ell=W_\ell D_{\ell-1}c^D_{\ell-1}+b_\ell$. The strict region $R_D=\{s:E^D_\ell(H^D_\ell s+c^D_\ell)>0\ \forall\ell\}$ is an open convex polyhedron, no longer a cone. \Cref{lem:consistency} holds with $F(s)=J_Ds+D_Lc^D_L$ on $R_D$: the induction is unchanged, and the derivative on $R_D$ is still $J_D$, because the biases do not enter the Jacobian. \Cref{lem:pl} is stated for linear pieces, but its proof uses only that $F$ agrees with an affine map $x\mapsto M_Cx+m_C$ on each cell: the telescoped differences $F(\gamma(u_j))-F(\gamma(u_{j-1}))=M_{C_j}(\gamma(u_j)-\gamma(u_{j-1}))$ do not see $m_C$. The nondegeneracy argument is unchanged, since each preactivation functional is affine with the same random linear part, and a dead formal piece now has a constant output and Jacobian zero. So \cref{eq:reduction} holds: $K_{n,L}(\alpha)=\max(\{0\}\cup\{\op{J_D}:R_D\ne\varnothing\})$.

\emph{Realizability.} With $c_D$ the stacked offsets, $R_D\ne\varnothing$ if and only if the affine subspace $\im(A_D)+c_D$ meets the orthant $O_D$. Under the gauge, extended to the biases by $b_\ell\mapsto\Sigma_\ell b_\ell$, one checks by the induction of \cref{lem:covariance} that $c^D_\ell(W^\Sigma,b^\Sigma)=\Sigma_\ell c^D_\ell(W,b)$, so the affine subspace transforms to $S_\Sigma(\im A_D+c_D)$; and the joint law of $(W,b)$ is preserved, because each bias coordinate is symmetric and independent of everything else. The proof of \cref{prop:orbit} then gives
\[
\E\bigl[F\one_{\{R_D\ne\varnothing\}}\bigr]=\E\Bigl[F\,\frac{r(\im A_D+c_D)}{2^{nL}}\Bigr]
\]
for every nonnegative gauge-invariant $F$, with $r$ counting the orthants met by the affine subspace.

\emph{Counting.} An affine subspace $A\subseteq\R^M$ of dimension $d\le n$ meets at most $\sum_{j=0}^d\binom Mj$ open orthants. Indeed the coordinate hyperplanes restrict to at most $M$ affine hyperplanes in $A$ (a coordinate that is constant on $A$ contributes none), the orthants met correspond to the regions of that arrangement as in \cref{sec:gauge}, and the number of regions of $M$ affine hyperplanes in dimension $d$ satisfies $R(M,d)\le R(M-1,d)+R(M-1,d-1)$ with $R(M,0)=1$ and $R(0,d)=1$, the same recurrence as \cref{eq:recur} with new boundary values, whose solution is $\sum_{j\le d}\binom Mj$ by Pascal's identity. With $M=nL$ and $d\le n$, \cref{lem:entropy} (with $k=n$, $N=nL$, $k/N=1/L\le\tfrac12$) gives $\sum_{j\le n}\binom{nL}j\le e^{nLH(1/L)}$, which is no larger than the bound $2e^{nLH(1/L)}$ used for $C_{n,L}$. For $L=1$ the count is $2^n$ as before.

\emph{The rest.} The net, the radial moments and the scalar bound concern $J_D$ only, so \cref{prop:master} and \cref{thm:finite} follow with the same budget.

\emph{The expansive side.} Homogeneity is lost, so \cref{eq:K-lower} is replaced by a ray argument. Let $F_0$ be the bias-free network with the same weights and fix $s\ne0$ at which every bias-free preactivation coordinate is nonzero (almost surely). For $c$ large enough the preactivations at the input $cs$ have the same signs as the bias-free preactivations at $s$, because $W_\ell(c\,x)+b_\ell=c\,(W_\ell x)+b_\ell$ and the bias is eventually dominated; so $cs$ lies in one region, on which $F$ is affine with linear part $J=DF$ equal to the bias-free Jacobian at $s$, and $F(cs)=cJs+\kappa$ with $\kappa$ independent of $c$. Since $Js=F_0(s)$ by homogeneity of $F_0$,
\[
K_{n,L}(\alpha)\ge\lim_{c\to\infty}\frac{\norm{F(cs)-F(0)}}{c\norm s}=\frac{\norm{F_0(s)}}{\norm s},
\]
and the proof of \cref{thm:above2} applies to $F_0$ unchanged.
\end{proof}

The biases play no role in the threshold because the Lipschitz constant is a statement about differences, and the only place they could enter, the realizability of a pattern, is handled by the same symmetry average with the same count. Deterministic or asymmetric biases are a different matter: the sign symmetry of the law is what the orbit identity uses, and without it the argument does not apply as written.

\section{One matrix, many steps: the finite-horizon norm law}
\label{sec:horizon}

The network of \cref{sec:intro} uses a fresh matrix at every layer. This section treats the other natural object: one Gaussian matrix applied again and again,
\begin{equation}\label{eq:fh-dyn}
x_{t+1}=\phi(Wx_t),\qquad t=0,1,2,\dots,\qquad W\in\R^{N\times N},\quad W_{ij}\overset{\text{ind}}{\sim}\mathcal N(0,1/N),
\end{equation}
started from a unit vector $x_0\in S^{N-1}$. In the notation of \cref{sec:intro} this is width $n=N$ and variance parameter $\alpha=1$, with the $L$ independent matrices replaced by a single one. The question is no longer the worst input but the ordinary one: what does the norm of the trajectory do? At $\alpha=1$ a layer halves the expected squared norm, and the theorem says that for a fixed number of steps this halving holds with exponentially small failure probability, for every deterministic input and for almost every input of one typical matrix. The proof is due to the same author as the global theorem and is given here in a condensed form; nothing is omitted from the logic, only from the commentary.

\subsection{Statements}

\begin{theorem}[deterministic input]\label{thm:fh}
Fix an integer $T\ge0$ and $\delta>0$. There are constants $C,c>0$, depending only on $T$ and $\delta$, such that for every $N\ge1$ and every deterministic $x_0\in S^{N-1}$,
\begin{equation}\label{eq:fh-main}
\PP_W\Bigl(\max_{0\le t\le T}\bigl|2^t\norm{x_t}^2-1\bigr|>\delta\Bigr)\le Ce^{-cN}.
\end{equation}
\end{theorem}

\begin{corollary}[random start independent of the matrix]\label{cor:fh-random}
If $X_0$ is a random unit vector independent of $W$, the same bound holds for $\PP_{W,X_0}$.
\end{corollary}

For a fixed matrix $W$ let $\mathcal B(W)=\{x_0\in S^{N-1}:\max_{t\le T}|2^t\norm{x_t}^2-1|>\delta\}$ be its set of bad starting points.

\begin{theorem}[almost every input]\label{thm:fh-ae}
Fix $T$ and $\delta$ and let $\mu_N$ be a probability measure on $S^{N-1}$, deterministic or random but independent of $W$. With the constants of \cref{thm:fh}, conditionally on $\mu_N$,
\[
\PP_W\bigl(\mu_N(\mathcal B(W))>e^{-cN/2}\bigr)\le Ce^{-cN/2}.
\]
\end{theorem}

The one-step case is the warm-up. For deterministic $x_0$, the coordinates of $Wx_0$ are independent $\mathcal N(0,1/N)$, so $\norm{x_1}^2=\frac1N\sum_i(\xi_i)_+^2$ with $\xi_i$ standard normal, and $\E(\xi_+)^2=\tfrac12\E\xi^2=\tfrac12$ by symmetry: the law of large numbers puts $\norm{x_1}^2$ near $\tfrac12$. The difficulty begins at the second step. The vector $x_1$ depends on $W$, so $Wx_1$ is not a Gaussian vector with independent coordinates, and the one-step argument cannot simply be repeated. This is the same entanglement as in \cref{sec:regions}, inside a single matrix. The proof handles it by conditioning, and the point of the section is why conditioning is legal here when it was avoided there.

\subsection{Tools}

Beyond \cref{sec:tools}, three estimates are needed. For $Z\sim\mathcal N(0,\sigma^2)$ and $t>0$, Markov's inequality on $e^{\lambda Z}$ with $\lambda=t/\sigma^2$ gives the Gaussian tail
\begin{equation}\label{eq:gtail}
\PP(\abs Z\ge t)\le2e^{-t^2/(2\sigma^2)}.
\end{equation}

\begin{lemma}[elementary Bernstein inequality]\label{lem:bernstein}
Let $X_1,\dots,X_N$ be independent, mean zero, with $\E\exp(\abs{X_i}/K)\le2$ for every $i$. Then for every $u>0$,
\[
\PP\Bigl(\Bigl|\frac1N\sum_{i=1}^NX_i\Bigr|>u\Bigr)\le2\exp\Bigl[-N\min\Bigl\{\frac{u^2}{16K^2},\frac u{8K}\Bigr\}\Bigr].
\]
\end{lemma}

\begin{proof}
Expanding the exponential, $2\ge\E e^{\abs{X_i}/K}=\sum_p\E\abs{X_i}^p/(p!K^p)$ termwise, so $\E\abs{X_i}^p\le2p!K^p$. For $\abs\lambda K\le\tfrac14$, since $\E X_i=0$,
\[
\E e^{\lambda X_i}=1+\sum_{p\ge2}\frac{\lambda^p\E X_i^p}{p!}\le1+2\sum_{p\ge2}(\abs\lambda K)^p=1+\frac{2(\abs\lambda K)^2}{1-\abs\lambda K}\le1+4\lambda^2K^2\le e^{4\lambda^2K^2}.
\]
By independence and Markov, $\PP(\frac1N\sum X_i>u)\le\exp(-\lambda Nu+4N\lambda^2K^2)$ for $0<\lambda\le1/(4K)$. If $u\le2K$ take $\lambda=u/(8K^2)$, which gives $e^{-Nu^2/(16K^2)}$; if $u>2K$ take $\lambda=1/(4K)$, which gives $\exp(-Nu/(4K)+N/4)\le e^{-Nu/(8K)}$. The lower tail is the same argument for $-X_i$.
\end{proof}

\begin{lemma}[average squared length of Gaussian vectors]\label{lem:chisq}
If $\Gamma_1,\dots,\Gamma_N$ are independent $\mathcal N(0,I_d)$ then $\PP\bigl(\frac1N\sum_i\norm{\Gamma_i}^2>2d\bigr)\le\exp(-\tfrac{1-\log2}2\,dN)$.
\end{lemma}

\begin{proof}
$S=\sum_i\norm{\Gamma_i}^2$ is a sum of $Nd$ independent squared standard normals, so $\E e^{S/4}=(1-\tfrac12)^{-Nd/2}=2^{Nd/2}$ by \cref{eq:mgf}, and $\PP(S>2Nd)\le e^{-Nd/2}2^{Nd/2}$.
\end{proof}

\subsection{What a query reveals: Gaussian innovations}

Write $w_i^T$ for the $i$-th row of $W$; the rows are independent $\mathcal N(0,I_N/N)$ vectors. A matrix--vector query $g=Wu$ with a unit vector $u$ returns $g_i=\langle w_i,u\rangle$, the coefficient of each row in the direction $u$ and nothing else: $w_i=\langle w_i,u\rangle u+P_{u^\perp}w_i$, and the query determines the first term only. After orthonormal queries $u_1,\dots,u_r$ the revealed part of every row is its projection onto $U_r=\mathrm{span}(u_1,\dots,u_r)$, and the unresolved residual is the projection onto $U_r^\perp$. The identity is pure geometry and holds when the directions are chosen adaptively. The probabilistic claim is that the residual is still Gaussian, and that needs care, because $U_r$ depends on $W$.

The mechanism in one line: if $A,B$ are independent standard normals and $s(A)\in\{\pm1\}$ is any measurable function of $A$ alone, then $C=s(A)B$ is standard normal and independent of $A$, because conditionally on $A$ the sign is fixed and $B$ is symmetric. The rule may depend on the past in any way; it may not look at the fresh variable it is applied to. The innovation theorem is the matrix version.

\begin{lemma}[one fresh direction inside an unobserved subspace]\label{lem:decomp}
Let $P$ be a deterministic orthogonal projection on $\R^N$, $u$ a unit vector with $Pu=0$, and $Z$ an $N\times N$ matrix with independent $\mathcal N(0,1/N)$ entries. Put $P^+=P+uu^T$, $\xi=Zu$ and $R=Z(I-P^+)$. Then $P^+$ is the orthogonal projection onto $\mathrm{range}(P)\oplus\mathrm{span}(u)$; $Z(I-P)=\xi u^T+R$; $\xi\sim\mathcal N(0,I_N/N)$; $\xi$ and $R$ are independent; and $R\eqd\widetilde Z(I-P^+)$ for an independent copy $\widetilde Z$ of $Z$.
\end{lemma}

\begin{proof}
$(P+uu^T)^2=P+uu^T$ because $Pu=0$ and $\norm u=1$, and $P+uu^T$ is symmetric, so it is the orthogonal projection onto the stated sum. Since $I-P=uu^T+(I-P^+)$, multiplying by $Z$ gives the identity. The $i$-th coordinate of $\xi$ is $z_i^Tu$ with $z_i$ the $i$-th row of $Z$, a centred Gaussian of variance $u^T(I_N/N)u=1/N$, independent across rows. The $i$-th row of $R$ is $r_i=(I-P^+)z_i$, jointly Gaussian with $\xi_i=z_i^Tu$, with cross-covariance $\E[r_i\xi_i]=(I-P^+)\E[z_iz_i^T]u=\frac1N(I-P^+)u=0$ because $P^+u=u$; jointly Gaussian and uncorrelated means independent, and the row pairs are independent across $i$, so $\xi$ is independent of the whole matrix $R$. Finally $R$ is a fixed linear image of $Z$, so replacing $Z$ by an independent copy does not change its law.
\end{proof}

Let $\mathcal F_0$ be trivial and $\mathcal F_k=\sigma(g_1,\dots,g_k)$; put $U_k=[u_1\cdots u_k]$, $P_k=U_kU_k^T$ and $M_k=\sum_{j\le k}g_ju_j^T$, so that $M_ku_a=g_a$ for $a\le k$: $M_k$ is the revealed action of $W$ on the queried span.

\begin{lemma}[Gaussian innovations under adaptive orthogonal queries]\label{lem:innovation}
Fix $K\le N$. Let $u_1\in S^{N-1}$ be deterministic and $g_1=Wu_1$. For $2\le k\le K$ let $u_k$ be $\mathcal F_{k-1}$-measurable with, almost surely, $\norm{u_k}=1$ and $\langle u_k,u_j\rangle=0$ for $j<k$, and put $g_k=Wu_k$. Then for every $k\le K$,
\begin{equation}\label{eq:strong}
W\mid\mathcal F_k\ \eqd\ M_k+\widetilde W_k(I-P_k),
\end{equation}
with $\widetilde W_k$ a fresh $\mathcal N(0,1/N)$ matrix independent of $\mathcal F_k$ (meaning $\E[\Phi(W)\mid\mathcal F_k]=\E_Z\Phi(M_k+Z(I-P_k))$ for bounded Borel $\Phi$, the revealed $M_k,P_k$ held fixed); $g_k\mid\mathcal F_{k-1}\sim\mathcal N(0,I_N/N)$; and $g_1,\dots,g_K$ are independent $\mathcal N(0,I_N/N)$ vectors.
\end{lemma}

\begin{proof}
Induction on $k$; at $k=0$ the statement is the law of $W$. Assume \cref{eq:strong} at $k-1$. Conditionally on $\mathcal F_{k-1}$, the vectors $g_1,\dots,g_{k-1}$, $u_1,\dots,u_k$ and the matrices $M_{k-1},P_{k-1}$ are fixed, $u_k$ because it is $\mathcal F_{k-1}$-measurable: it may have been chosen by an elaborate rule, but once the past is revealed the rule has produced a known vector. Multiply \cref{eq:strong} by $u_k$: $M_{k-1}u_k=\sum_{j<k}g_j\langle u_j,u_k\rangle=0$ by orthogonality, and $(I-P_{k-1})u_k=u_k$, so $g_k=Wu_k\eqd Zu_k$ with $Z$ fresh; conditionally on the past $u_k$ is a fixed unit vector, so $g_k\mid\mathcal F_{k-1}\sim\mathcal N(0,I_N/N)$, a law that does not depend on the past: $g_k$ is independent of $\mathcal F_{k-1}$. For the representation at step $k$ apply \cref{lem:decomp} with $P=P_{k-1}$ and $u=u_k$: $Z(I-P_{k-1})=\xi u_k^T+R$ with $\xi=Zu_k=g_k$ and $R\eqd\widetilde Z(I-P_k)$ independent of $\xi$. So, conditionally on $\mathcal F_{k-1}$, the pair $(g_k,W)$ has the law of $(\xi,M_{k-1}+\xi u_k^T+R)$ with $\xi$ independent of $R$; conditioning further on $\xi=g_k$ leaves the law of $R$ unchanged and gives $W\mid\mathcal F_k\eqd M_{k-1}+g_ku_k^T+\widetilde Z(I-P_k)=M_k+\widetilde Z(I-P_k)$. Joint independence: for bounded $\varphi_j$, $\E\prod_{j\le K}\varphi_j(g_j)=\E[\prod_{j<K}\varphi_j(g_j)\,\E[\varphi_K(g_K)\mid\mathcal F_{K-1}]]=\E[\prod_{j<K}\varphi_j(g_j)]\,\E\varphi_K(G)$ with $G\sim\mathcal N(0,I_N/N)$, and repeat.
\end{proof}

Three hypotheses, each necessary. Orthogonality: if $u_2=u_1$ then $g_2=g_1$; in general $Wu_k=\sum_{j<k}\alpha_jg_j+Wu_k^\perp$ and only the orthogonal part is fresh, which is what Gram--Schmidt isolates. The direction may depend only on the past: a rule allowed to inspect the unobserved residual could pick the direction of its largest projection. Gaussianity: for a non-Gaussian isotropic vector, orthogonal projections can be uncorrelated without being independent. And the matrix model matters: for a symmetric Gaussian matrix the rows are coupled and the statement is false as written.

\subsection{Gram--Schmidt on the trajectory}

Let $V_t=\mathrm{span}(x_0,\dots,x_t)$, of dimension at most $t+1$. Apply Gram--Schmidt to $x_0,x_1,\dots$ in order: $v_1=x_0$; when a new $x_t$ has a nonzero residual $\rho_t=x_t-\sum_{j\le r}\langle x_t,v_j\rangle v_j$, set $v_{r+1}=\rho_t/\norm{\rho_t}$. Each new direction is a unit vector orthogonal to the earlier ones, and it is a measurable function of the images $g_j=Wv_j$ already exposed: $x_1=\phi(g_1)$, and inductively $Wx_t=\sum_j\langle x_t,v_j\rangle g_j$ so $x_{t+1}=\phi(Wx_t)$, its projection onto $V_t$, its residual and the normalized residual are all functions of $g_1,\dots,g_r$. So the Gram--Schmidt directions are admissible probes for \cref{lem:innovation}.

One technical point keeps the probe count deterministic. \Cref{lem:innovation} is stated for a fixed number $K$ of probes; the number of directions the trajectory actually creates is random. So the list is padded: with $d=T+1\le N$, after the trajectory's directions are exhausted, further probes are chosen by a fixed rule from the transcript (the first standard basis vector not yet in the span, projected and normalized). This produces exactly $d$ transcript-measurable orthonormal probes $v_1,\dots,v_d$, with $V_T\subseteq\mathrm{span}(v_1,\dots,v_d)$ and every padding vector orthogonal to $V_T$; by \cref{lem:innovation}, $g_1,\dots,g_d$ are independent $\mathcal N(0,I_N/N)$. For $0\le t\le T$ define the coefficient vector $a_t=(\langle x_t,v_1\rangle,\dots,\langle x_t,v_d\rangle)\in\R^d$; the padding coordinates are zero, Parseval gives $\norm{a_t}=\norm{x_t}$, and the driving representation is
\begin{equation}\label{eq:driving}
Wx_t=\sum_{j=1}^d(a_t)_jg_j .
\end{equation}
For each output coordinate $i$ let $\Gamma_i=\sqrt N\bigl((g_1)_i,\dots,(g_d)_i\bigr)\in\R^d$; the $Nd$ scalars $\sqrt N(g_j)_i$ are independent standard normals, so $\Gamma_1,\dots,\Gamma_N$ are independent $\mathcal N(0,I_d)$. Taking the $i$-th coordinate of \cref{eq:driving}, applying ReLU and homogeneity, squaring and summing,
\begin{equation}\label{eq:onestep}
\norm{x_{t+1}}^2=\frac1N\sum_{i=1}^N(a_t\cdot\Gamma_i)_+^2 .
\end{equation}
The whole trajectory is now written in terms of $d$ independent Gaussian vectors and $T+1$ coefficient vectors that live in a space of dimension $d=T+1$, a dimension that does not grow with $N$. The coefficient vectors are random and depend on the same Gaussians; a concentration estimate for one fixed $a$ would not be enough, and the next lemma is uniform over a ball.

\subsection{A uniform law of large numbers in a low dimension}

\begin{lemma}[uniform Gaussian law of large numbers]\label{lem:ulln}
Fix $d$, $R<\infty$ and $\varepsilon>0$, and let $\Gamma_1,\dots,\Gamma_N$ be independent $\mathcal N(0,I_d)$. There are $C,c>0$ depending only on $d,R,\varepsilon$ such that
\[
\PP\Bigl(\sup_{\norm a\le R}\Bigl|\frac1N\sum_{i=1}^N(a\cdot\Gamma_i)_+^2-\frac12\norm a^2\Bigr|>\varepsilon\Bigr)\le Ce^{-cN}.
\]
\end{lemma}

\begin{proof}
Write $A_N(a)=\frac1N\sum_i(a\cdot\Gamma_i)_+^2$. For fixed $a$, $Z=a\cdot\Gamma$ is $\mathcal N(0,\norm a^2)$ and $\E Z_+^2=\tfrac12\E Z^2=\tfrac12\norm a^2$ by symmetry, so $\E A_N(a)=\tfrac12\norm a^2$. Concentration for one $a$ with $\norm a\le R$: $Y_i=(a\cdot\Gamma_i)_+^2$ satisfies $0\le Y_i\le Z_i^2$ and $\E Y_i\le R^2/2$, so $X_i=Y_i-\E Y_i$ has $\abs{X_i}\le Z_i^2+R^2/2$, and with $K_R=32R^2$, by \cref{eq:mgf}, $\E e^{\abs{X_i}/K_R}\le e^{1/64}(1-2R^2/(32R^2))^{-1/2}=e^{1/64}(1-\tfrac1{16})^{-1/2}<2$. \Cref{lem:bernstein} with $u=\varepsilon/2$ gives $\PP(\abs{A_N(a)-\tfrac12\norm a^2}>\varepsilon/2)\le2e^{-c_1N}$ with $c_1=\min\{\varepsilon^2/(64K_R^2),\varepsilon/(16K_R)\}$, which depends on $R,\varepsilon$ only. Nearby coefficient vectors: for $\norm a,\norm b\le R$ and any $\Gamma$, using $\abs{u^2-v^2}=\abs{u-v}\abs{u+v}$, the $1$-Lipschitz property of $\phi$ and Cauchy--Schwarz, $\abs{(a\cdot\Gamma)_+^2-(b\cdot\Gamma)_+^2}\le2R\norm{a-b}\norm\Gamma^2$; averaging,
\[
\abs{A_N(a)-A_N(b)}\le2R\norm{a-b}\,\frac1N\sum_i\norm{\Gamma_i}^2,
\]
and the means differ by at most $R\norm{a-b}$. Net: put $\alpha_0=\min\{R,\,\varepsilon/(2R(4d+1))\}$. The ball $B_d(R)$ has an $\alpha_0$-net $\net$ of size $M\le(3R/\alpha_0)^d$ (a maximal $\alpha_0$-separated set, by the volume argument of \cref{lem:net} with radii $\alpha_0/2$ and $R+\alpha_0/2$), a number independent of $N$. On the event that every net point concentrates to within $\varepsilon/2$, which fails with probability at most $2Me^{-c_1N}$, and the event $\frac1N\sum\norm{\Gamma_i}^2\le2d$ of \cref{lem:chisq}, which fails with probability at most $e^{-c_2N}$, take any $a$ and a net point $b$ within $\alpha_0$: $\abs{A_N(a)-\tfrac12\norm a^2}\le4Rd\alpha_0+\varepsilon/2+R\alpha_0=R(4d+1)\alpha_0+\varepsilon/2\le\varepsilon$. So the bound holds with $C=2M+1$ and $c=\min\{c_1,c_2\}$.
\end{proof}

\subsection{A crude operator-norm bound}

\begin{lemma}\label{lem:opnorm}
$\PP(\op W>8)\le2\exp(-(8-\log81)N)$.
\end{lemma}

\begin{proof}
If $\mathcal M$ is an $\epsilon$-net of $S^{N-1}$ with $\epsilon<\tfrac12$, then $\op A\le(1-2\epsilon)^{-1}\max_{u,v\in\mathcal M}\abs{u^TAv}$: for unit $x,y$ and net points $u,v$ within $\epsilon$, $\abs{y^TAx}\le\abs{v^TAu}+\abs{(y-v)^TAx}+\abs{v^TA(x-u)}\le\max\abs{p^TAq}+2\epsilon\op A$, and take the supremum. A $\tfrac14$-net has at most $9^N$ points by the volume argument. For fixed unit $u,v$, $u^TWv=\sum_{ij}u_iW_{ij}v_j$ is centred Gaussian of variance $\frac1N\sum_iu_i^2\sum_jv_j^2=1/N$, so $\PP(\abs{u^TWv}>4)\le2e^{-8N}$ by \cref{eq:gtail}; a union bound over at most $81^N$ pairs gives $\PP(\op W>8)\le81^N\cdot2e^{-8N}$, and $8-\log81>0$.
\end{proof}

\subsection{Proof of the deterministic theorem}

\begin{proof}[Proof of \cref{thm:fh}]
For $T=0$ the bad event is empty. Let $T\ge1$ and first $N\ge T+1$. On the event $\mathcal O=\{\op W\le8\}$, whose complement has probability at most $C_{\rm op}e^{-c_{\rm op}N}$ by \cref{lem:opnorm}, $\norm{x_{t+1}}=\norm{\phi(Wx_t)}\le\norm{Wx_t}\le8\norm{x_t}$, so $\norm{x_t}\le8^t\le8^T=:R$ for all $t\le T$. Build the $d=T+1$ probes and innovations of the previous subsection; on $\mathcal O$ every coefficient vector satisfies $\norm{a_t}=\norm{x_t}\le R$. Put $\eta=\delta2^{-(T+2)}$ and apply \cref{lem:ulln} with $d=T+1$, radius $R=8^T$ and tolerance $\eta$: outside an event $\mathcal E^c$ of probability at most $C_{\rm LLN}e^{-c_{\rm LLN}N}$, the inequality $\abs{A_N(a)-\tfrac12\norm a^2}\le\eta$ holds for every $a$ in the ball at once. The coefficient vector $a_t$ is random and depends on the Gaussians, and this is exactly why uniformity is needed: on $\mathcal O\cap\mathcal E$ the realized $a_t$ lies in the ball, so it may be substituted, and \cref{eq:onestep} gives
\[
\bigl|\norm{x_{t+1}}^2-\tfrac12\norm{x_t}^2\bigr|\le\eta,\qquad0\le t\le T-1 .
\]
Write $q_t=\norm{x_t}^2$ and $e_t=\abs{q_t-2^{-t}}$; then $e_0=0$ and $e_{t+1}\le\abs{q_{t+1}-\tfrac12q_t}+\tfrac12\abs{q_t-2^{-t}}\le\eta+\tfrac12e_t$, so $e_t\le\eta(1+\tfrac12+\dots)\le2\eta$, and $\abs{2^tq_t-1}=2^te_t\le2^{T+1}\eta=\delta/2<\delta$ for every $t\le T$. The bad event is therefore contained in $\mathcal O^c\cup\mathcal E^c$, of probability at most $Ce^{-cN}$ with constants depending on $T,\delta$ only. For the finitely many $N<T+1$ the probability is at most $1$, and enlarging $C$ makes $Ce^{-cN}\ge1$ there.
\end{proof}

\subsection{A random start, and almost every input}

\begin{proof}[Proof of \cref{cor:fh-random}]
Condition on $X_0=x$. Independence leaves the law of $W$ unchanged, so \cref{eq:fh-main} holds with the same constants for every $x$, and taking the expectation over $X_0$ gives the bound.
\end{proof}

\begin{proof}[Proof of \cref{thm:fh-ae}]
Let $B(x_0,W)$ be the indicator that $x_0$ is bad for $W$. It is measurable: $(x_0,W)\mapsto x_t$ is continuous for each $t$, so $\max_t\abs{2^t\norm{x_t}^2-1}$ is continuous and its superlevel set is Borel. For deterministic $\mu_N$ put $M(W)=\int B(x_0,W)\,\mu_N(dx_0)=\mu_N(\mathcal B(W))\in[0,1]$. By Tonelli and the uniform bound \cref{eq:fh-main},
\[
\E_WM(W)=\int\PP_W\bigl(B(x_0,W)=1\bigr)\,\mu_N(dx_0)\le Ce^{-cN},
\]
and Markov's inequality with threshold $e^{-cN/2}$ gives $\PP_W(M(W)>e^{-cN/2})\le e^{cN/2}\cdot Ce^{-cN}=Ce^{-cN/2}$. For a random $\mu_N$ independent of $W$, condition on its realized value: the law of $W$ is unchanged and the deterministic case applies with the same constants.
\end{proof}

\subsection{With a bias vector}

Let $x_{t+1}=\phi(Wx_t+b)$, with one bias vector reused at every step, independent of $W$, and $b_i\overset{\mathrm{ind}}{\sim}\mathcal N(0,\sigma^2/N)$. The reference profile is
\[
\bar q_0=1,\qquad\bar q_{t+1}=\tfrac12(\bar q_t+\sigma^2),\qquad\text{so}\qquad\bar q_t=\sigma^2+(1-\sigma^2)2^{-t},
\]
whose fixed point is $\sigma^2$. The result below controls a fixed horizon; it makes no assertion about the infinite-time trajectory.

\begin{theorem}[finite horizon, with a bias]\label{thm:fh-bias}
Fix an integer $T\ge0$, $\delta>0$ and $\sigma\ge0$. There are $C,c>0$ depending only on $T,\delta,\sigma$ such that for every $N\ge1$ and every deterministic $x_0\in S^{N-1}$,
\[
\PP_{W,b}\Bigl(\max_{0\le t\le T}\bigl|\norm{x_t}^2-\bar q_t\bigr|>\delta\Bigr)\le Ce^{-cN},
\]
The analogues of \cref{cor:fh-random,thm:fh-ae} hold for a random start or input measure independent of the pair $(W,b)$, with bad inputs defined by the displayed event.
\end{theorem}

\begin{proof}
The case $\sigma=0$ follows from \cref{thm:fh}, and $T=0$ is immediate. Suppose $\sigma>0$, $T\ge1$ and $N\ge T+1$. Apply the innovation argument conditionally on $b$, taking $\sigma(b)$ as the initial information. Since $W$ is independent of $b$, the conditional law of the $d=T+1$ padded innovations is the same product Gaussian law for every $b$. Consequently those innovations are jointly independent of $b$ after removing the conditioning.

Put $\beta_i=\sqrt N\,b_i$, $\Gamma_i'=(\Gamma_i,\beta_i/\sigma)$ and $a_t'=(a_t,\sigma)$. The vectors $\Gamma_i'$ are therefore independent $\mathcal N(0,I_{d+1})$, and the driving identity becomes
\[
\norm{x_{t+1}}^2=\frac1N\sum_{i=1}^N(a_t'\cdot\Gamma_i')_+^2,
\qquad \norm{a_t'}^2=\norm{x_t}^2+\sigma^2.
\]
On $\{\op W\le8,\ \norm b^2\le2\sigma^2\}$, the recursion $\norm{x_{t+1}}\le8\norm{x_t}+2\sigma$ gives $\norm{a_t'}\le R':=8^T(1+3\sigma)$ for $t\le T$. The complement has probability at most $Ce^{-cN}$ by \cref{lem:opnorm,lem:chisq}. Apply \cref{lem:ulln} in dimension $d+1$ with radius $R'$ and tolerance $\eta=\delta/4$. Its uniformity permits the random coefficient $a_t'$ to depend on all the Gaussian vectors. Outside another event of probability at most $Ce^{-cN}$,
\[
\left|\norm{x_{t+1}}^2-\tfrac12(\norm{x_t}^2+\sigma^2)\right|\le\eta
\quad(0\le t<T).
\]
Hence $e_t:=|\norm{x_t}^2-\bar q_t|$ satisfies $e_0=0$ and $e_{t+1}\le\eta+e_t/2$, so $e_t\le2\eta<\delta$. Enlarge $C$ for the finitely many $N<T+1$. Finally, integrate the deterministic-input bound over a start or measure independent of $(W,b)$ and apply Tonelli and Markov exactly as before, now over the joint randomness $(W,b)$.
\end{proof}

\subsection{What the theorem says and does not say}

The factor $\tfrac12$ is $\E(\xi_+)^2=\tfrac12\E\xi^2$ for a symmetric $\xi$: ReLU keeps half of the expected squared magnitude, and the theorem says that this average concentrates along the adaptively chosen trajectory. The horizon is fixed because the trajectory spans a subspace of dimension at most $T+1$, the uniform law of large numbers is applied in that dimension, and the size of the net is then a constant independent of $N$; a horizon growing with $N$ would make the net and the constants grow, and the present argument would need more. The deterministic theorem is uniform in the choice of one starting vector; it does not say that one matrix works for every point of the sphere at once, and \cref{thm:fh-ae} is the weaker simultaneous statement it does give: for one typical matrix, all but an exponentially small set of inputs are good. The start, or the measure, may be conditioned on only because it is independent of $W$; a start chosen after seeing the matrix could sit on an exceptional direction.

\section{How the two results fit together}
\label{sec:together}

\subsection{The same object, two questions}

Both results are about a Gaussian matrix with independent $\mathcal N(0,\alpha/n)$ entries followed by ReLU, and both live at the scale where a layer keeps a squared norm of order one. Both rest on the same elementary fact, $\E(\xi_+)^2=\tfrac12\E\xi^2$ for a symmetric $\xi$: ReLU discards half of the expected squared magnitude, so a layer multiplies a typical squared norm by $\alpha/2$. In \cref{sec:horizon} this is the number $\tfrac12$ at $\alpha=1$; in \cref{sec:theorems} it is the reason the critical variance is $2$. Both face the same obstacle. The input to a layer is a function of the matrices before it, so the vector a layer acts on is not independent of the layer; in the tied case the vector is a function of the very matrix that acts on it. Neither proof pretends otherwise. Both replace an infinite supremum by a finite net and pay a fixed exponential price for it, both produce failure probabilities exponentially small in the width, and both fix the number of layers first and let the width grow, with the horizon $T$ or the depth $L$ entering the constants.

The questions are different, and the difference decides the method.

\begin{center}
\begin{tabular}{p{0.28\textwidth}p{0.32\textwidth}p{0.32\textwidth}}
\toprule
 & global stability (\cref{thm:finite,thm:main}) & finite horizon (\cref{thm:fh,thm:fh-ae})\\
\midrule
matrices & $L$ independent, one per layer & one matrix, applied $T$ times\\
variance & $\alpha/n$, any $\alpha$ & $1/N$, that is $\alpha=1$\\
quantified over & every input at once (a supremum) & one input; then almost every input\\
quantity controlled & the Lipschitz constant, a derivative & the norm of the trajectory\\
conclusion & $K_{n,L}(\alpha)\le1$ w.h.p. for $\alpha<2$, $L\ge L_0(\alpha)$ & $\norm{x_t}^2\approx2^{-t}$ for $t\le T$\\
the obstacle & the worst input's masks depend on the weights & $x_t$ depends on $W$\\
the treatment & never condition: average the realizability event over the sign gauges & condition exactly: Gram--Schmidt probes give independent innovations\\
the counting & orthants met by an $n$-dimensional subspace of $\R^{nL}$, $\le e^{nLH(1/L)}$ & a net of a ball in $\R^{T+1}$, size independent of $N$\\
the moment & one exact radial moment per layer, order $2\theta n$ & a Bernstein bound for each net point\\
failure probability & $2e^{-c_L(\alpha)n}$ beyond $L_0(\alpha)$ & $Ce^{-cN}$ for every $N$\\
\bottomrule
\end{tabular}
\end{center}

\subsection{Why each method fits its question}

The finite-horizon proof conditions on what the matrix has revealed, and it may do so because the revealed information is finite-dimensional: the trajectory queries $W$ in at most $T+1$ directions, Gram--Schmidt makes them orthogonal, and the innovation lemma says that orthogonal queries chosen from the past return independent Gaussian answers. The dimension of the queried subspace is the horizon, not the width, and the uniform law of large numbers is applied in that small dimension. The argument cannot be run over all inputs at once, because the set of all inputs queries $W$ in every direction, and no fixed-length transcript describes it. This is the reason the deterministic theorem is about one input, and why the extension to almost every input comes through Fubini, from a typical matrix's point of view, rather than through a supremum.

The global theorem asks about the worst input, so it cannot condition on a trajectory; there is no trajectory, there are $2^{nL}$ candidate activation patterns and the adversary picks among them after seeing the weights. The gauge argument refuses to condition at all. It fixes a pattern on paper, uses a symmetry of the law of the weights to average the realizability event over all $2^{nL}$ orthants, and turns the average into a deterministic count. The price is the count, $e^{nLH(1/L)}$ against the $2^{nL}$ patterns, and the argument works only when the radial contraction, one exact moment per independent layer, outweighs it, which is what the budget $B_L$ measures and why a depth $L_0(\alpha)$ appears. The symmetry needs independent layers: the gauge inserts a sign matrix on the left of one layer and cancels it on the right of the next, and with a single tied matrix the available symmetry is conjugation $W\mapsto\Sigma W\Sigma$ by one sign matrix, whose orbit does not move a target orthant through all orthants. So the two methods are not interchangeable. One is a conditioning argument that works along a low-dimensional trajectory for one input; the other is a symmetry argument that works over all inputs for independent layers. What is quantified over decides which one applies.

\subsection{Where the second method sharpens the first theorem}

The expansive side of the global theorem, \cref{thm:above2}, is a one-input statement: at $\alpha>2$ a fixed input's forward norm grows, so the Lipschitz constant exceeds one. It was proved with Chebyshev's inequality, which gives a $1/n$ rate. The finite-horizon proof handles the same one-input forward pass with exponential concentration, and its ingredient, the moment-generating function of a masked Gaussian square, is already in \cref{lem:mgf}. Using it in the global setting gives an exponential rate.

\begin{corollary}[the expansive side, exponentially]\label{cor:chernoff}
For $\alpha>2$ let $c=1/\alpha<\tfrac12$ and
\[
I(c)=\sup_{s>0}\Bigl[-sc-\log\frac{1+(1+2s)^{-1/2}}2\Bigr]>0 .
\]
Then for every $n\ge1$ and $L\ge1$,
\[
\PP\bigl(K_{n,L}(\alpha)\le1\bigr)\le L\,e^{-nI(1/\alpha)} .
\]
\end{corollary}

\begin{proof}
As in the proof of \cref{thm:above2}, $\PP(K\le1)\le L\,\PP(S_n\le n/\alpha)$ with $S_n=\sum_iB_iZ_i^2$. For $s>0$, Markov's inequality on $e^{-sS_n}$ and \cref{lem:mgf} with $t=-s$ give
\[
\PP(S_n\le nc)=\PP\bigl(e^{-sS_n}\ge e^{-snc}\bigr)\le e^{snc}\Bigl[\frac{1+(1+2s)^{-1/2}}2\Bigr]^n,
\]
and optimizing over $s$ gives $e^{-nI(c)}$. Positivity: $\log\frac{1+(1+2s)^{-1/2}}2=-\tfrac s2+O(s^2)$ as $s\downarrow0$, so the bracket is $s(\tfrac12-c)+O(s^2)>0$ for small $s$ when $c<\tfrac12$.
\end{proof}

The rate $I(1/\alpha)$ is a number: \cref{sec:numerics} lists it for several $\alpha$ and compares the two bounds. In the tied-weight dynamics the corresponding statement is \cref{thm:fh} itself, and its constants come from the same Bernstein and Chernoff estimates; the only difference is the innovation step that makes the one-input pass legitimate when the matrix is reused.

\subsection{What neither result says}

The global theorem needs independent layers; the finite-horizon theorem needs one input and a fixed horizon. Whether the tied dynamics $x\mapsto\phi(Wx)$ is non-expansive as a map, that is, whether a statement like \cref{thm:main} holds for the global Lipschitz constant of $\phi(W\cdot)^{\circ L}$, is not settled by either argument: the gauge orbit of one matrix is too small, and a supremum over inputs is not a finite transcript. That question is open here, and it is the natural meeting point of the two methods.

\section{Numerics}
\label{sec:numerics}

Every number in this section was computed for this paper: the exact quantities by evaluating the formulas, the Monte Carlo quantities by simulating the network with a fixed seed and comparing with the formulas they are supposed to obey. Where a quantity was sampled rather than enumerated the table says so, since a sampled supremum is a lower bound.

\subsection{The orthant count}

\Cref{tab:cover} draws a random subspace $U$ of dimension $d$ in $\R^M$ and counts the open orthants it meets. In the plane the count is exact, by enumerating the arcs between the angles at which the coordinate functionals vanish; for $d\ge3$ the count is a lower bound from sampled directions. The bound of \cref{lem:cover} is attained whenever the count is exact, in agreement with the general-position remark.

\begin{table}[h]\centering
\begin{tabular}{rrrrrl}
\toprule
$M$ & $d$ & $2^M$ & $2\sum_{j<d}\binom{M-1}j$ & observed $r(U)$, mean (max) & method\\
\midrule
6 & 2 & 64 & 12 & 12.0 (12) & exact\\
8 & 2 & 256 & 16 & 16.0 (16) & exact\\
10 & 2 & 1024 & 20 & 20.0 (20) & exact\\
20 & 2 & 1048576 & 40 & 40.0 (40) & exact\\
6 & 3 & 64 & 32 & 31.9 (32) & sampled\\
9 & 3 & 512 & 74 & 73.4 (74) & sampled\\
12 & 3 & 4096 & 134 & 131.8 (134) & sampled\\
8 & 4 & 256 & 128 & 120.8 (127) & sampled\\
\bottomrule
\end{tabular}
\caption{Orthants met by a random $d$-dimensional subspace of $\R^M$ against the bound of \cref{lem:cover}; 30 subspaces per row for $d=2$, 12 for $d\ge3$.}
\label{tab:cover}
\end{table}

For the network itself, the identity \cref{eq:orbit} with $F=1$, summed over all masks, says that the expected number of nonempty strict regions is $\sum_D\E\,r(\im A_D)/2^{nL}\le C_{n,L}$. \Cref{tab:regions} compares the number of realized masks with $C_{n,L}$ and with the number $2^{nL}$ of formal masks. Even at width two the realized count stays far below the formal count, and below $C_{n,L}$.

\begin{table}[h]\centering
\begin{tabular}{rrrrrl}
\toprule
$n$ & $L$ & formal masks $2^{nL}$ & $C_{n,L}$ & realized masks, mean (min--max) & method\\
\midrule
2 & 1 & 4 & 4 & 4.0 (4--4) & exact\\
2 & 2 & 16 & 8 & 4.9 (4--6) & exact\\
2 & 3 & 64 & 12 & 5.0 (4--7) & exact\\
2 & 4 & 256 & 16 & 5.5 (4--8) & exact\\
2 & 5 & 1024 & 20 & 5.3 (4--8) & exact\\
2 & 6 & 4096 & 24 & 5.5 (4--10) & exact\\
3 & 2 & 64 & 32 & 15.5 (8--20) & sampled\\
3 & 3 & 512 & 74 & 20.3 (11--28) & sampled\\
4 & 2 & 256 & 128 & 47.1 (24--71) & sampled\\
\bottomrule
\end{tabular}
\caption{Nonempty strict regions of the random network (60 draws per row at $n=2$, 25 otherwise).}
\label{tab:regions}
\end{table}

At fixed $L$ the per-width exponent of $C_{n,L}$ converges to $LH(1/L)$ (\cref{lem:C-rate}). At $n=40$: for $L=3$, $\frac1n\log C_{n,L}=1.851$ against $LH(1/L)=1.910$ and $L\log2=2.079$; for $L=10$, $3.146$ against $3.251$ and $6.931$; for $L=100$, $5.434$ against $5.600$ and $69.31$. The entropy bound of \cref{lem:C-entropy} is within a few percent of the exact count per unit width already at width 40, and the saving against the $2^{nL}$ formal masks is the whole difference between $\log L+1$ and $L\log2$.

\subsection{The realizable-mask identity}

\Cref{tab:orbit} is the paired test of \cref{prop:orbit} at width $n=2$ described in \cref{sec:gauge}: for a fixed mask $D$, the same weight draws are used on both sides, $r(\im A_D)$ is computed exactly by arc enumeration, and the paired difference is compared with its standard error. Two functionals are used, $F\equiv1$ and $F=\min(\op{J_D},1)$, both bounded, so the standard error is meaningful.

\begin{table}[h]\centering
\begin{tabular}{rlrrrrr}
\toprule
$L$ & $F$ & draws & $\E[F\one_{\{R_D\ne\varnothing\}}]$ & $\E[F\,r(\im A_D)/2^{nL}]$ & ratio & paired $z$\\
\midrule
2 & $1$ & 200000 & 0.24878 & 0.25000 & 0.995 & $-1.27$\\
3 & $1$ & 200000 & 0.06283 & 0.06250 & 1.005 & $+0.59$\\
4 & $1$ & 150000 & 0.01589 & 0.01563 & 1.017 & $+0.83$\\
5 & $1$ & 120000 & 0.00390 & 0.00391 & 0.998 & $-0.04$\\
3 & $\min(\op{J_D},1)$ & 200000 & 0.02307 & 0.02296 & 1.005 & $+0.41$\\
4 & $\min(\op{J_D},1)$ & 150000 & 0.00492 & 0.00493 & 0.999 & $-0.03$\\
\bottomrule
\end{tabular}
\caption{The exact orbit identity \cref{eq:orbit} at width 2, the mask $D$ fixed, the subspace $\im A_D$ enumerated exactly.}
\label{tab:orbit}
\end{table}

Two more checks that the proof leans on directly were run alongside: summing the indicator $\one_{\{S_\Sigma U\cap O_D\ne\varnothing\}}$ over the whole gauge group returned exactly $r(U)$ in every draw, and $\op{J_D}$ was unchanged by a random gauge in every draw, to machine precision.

The all-active formula \cref{eq:all-active} gives $\PP(R_D\ne\varnothing)=\tfrac12$ at $(n,L)=(2,2)$ and $(3,2)$, $0.1445$ at $(3,3)$, $0.0898$ at $(5,3)$, $0.0307$ at $(10,3)$, $0.00432$ at $(20,3)$ and $3.0\times10^{-18}$ at $(10,10)$: the chance that a deep random network has an input at which every unit fires at every layer is exactly computable and small.

\subsection{The moment formulas}

\Cref{tab:radial} checks the exact radial factor \cref{eq:radial} and the mixture identity \cref{eq:mixture} by brute force.

\begin{table}[h]\centering\small
\begin{tabular}{llrrr}
\toprule
identity & parameters & Monte Carlo & exact & ratio\\
\midrule
\cref{eq:radial} & $n=4$, $L=3$, $p=1$, $\alpha=1.3$, masks $1101,1001,1110$ & 0.6247 & 0.6179 & 1.011\\
\cref{eq:radial} & $n=4$, $L=3$, $p=2$, same masks & 2.063 & 2.121 & 0.973\\
\cref{eq:radial} & $n=6$, $L=2$, $p=1.5$, $\alpha=1.3$, masks $101101,011011$ & 0.9022 & 0.8987 & 1.004\\
\cref{eq:mixture} & $n=6$, $p=1$, $\alpha=1.3$ & 41.65 & 41.60 & 1.001\\
\cref{eq:mixture} & $n=6$, $p=2.5$, $\alpha=1.3$ & 61.99 & 61.93 & 1.001\\
\cref{eq:mixture} & $n=12$, $p=1.7$, $\alpha=2$ & 5087 & 5079 & 1.002\\
\bottomrule
\end{tabular}
\caption{The fixed-vector radial moments and the scalar mixture identity, $2\times10^5$ draws each; the masks are the diagonals of $D_1,\dots,D_L$; for \cref{eq:mixture} both sides are multiplied by $2^n$.}
\label{tab:radial}
\end{table}

The master inequality \cref{eq:master} is far from tight at small width, and the table says by how much: at $(n,L,\alpha,p)=(2,2,1,1)$ the left side $\E K^{2p}$ is $1.51$ against a bound of $200$; at $(2,3,1,1.5)$, $1.94$ against $557$; at $(2,2,0.6,2)$, $0.88$ against $318$. The slack is the product of the net factor $5^n4^p$ and of $C_{n,L}$ counting every orthant a subspace could meet rather than the ones it does. The inequality is designed for the exponential scale, where these factors cost a fixed rate per unit width, and it is there that it is used.

The $\tfrac12$-net of \cref{lem:net} is also loose by a fixed exponential factor: a greedy maximal $\tfrac12$-separated set on $S^{n-1}$ had $2,9,34,102,279,676$ points for $n=1,\dots,6$ against the bound $5^n$; the ratio behaves like $e^{-0.5n}$. Neither slack affects the threshold, which is decided by the $L$-fold radial factor.

\subsection{The finite-width bound}

\Cref{tab:threshold} lists the threshold $\alpha^*_L$ of \cref{eq:alpha-star}, the largest $\alpha$ at which the exact budget $B_L$ is negative for some $\theta$, together with the asymptotic formula of \cref{thm:main}(iv) evaluated at the same $L$. The asymptotic formula overstates the bound at moderate depth: at $L=100$ it reads $1.01$ with the $o(1)$ dropped, or $0.81$ with the constant $A=1+\log5+1$ kept, while the budget delivers $0.68$. At $L=10^5$ the three agree to a few parts in a thousand.

\begin{table}[h]\centering
\begin{tabular}{rrrr}
\toprule
$L$ & $\alpha^*_L$ (exact budget) & $2e^{-\sqrt{10\log L/L}}$ & $2e^{-\sqrt{10(\log L+A)/L}}$, $A=2+\log5$\\
\midrule
10 & 0.0412 & 0.4386 & 0.1758\\
30 & 0.2502 & 0.6896 & 0.4336\\
100 & 0.6778 & 1.0146 & 0.8080\\
300 & 1.0822 & 1.2932 & 1.1457\\
1000 & 1.4261 & 1.5378 & 1.4461\\
3000 & 1.6390 & 1.6986 & 1.6428\\
$10^4$ & 1.7875 & 1.8170 & 1.7859\\
$10^5$ & 1.9254 & 1.9333 & 1.9237\\
$10^6$ & 1.9744 & 1.9766 & 1.9738\\
\bottomrule
\end{tabular}
\caption{The threshold the finite theorem delivers at depth $L$, against the large-$L$ formula.}
\label{tab:threshold}
\end{table}

The other reading of the same budget is the depth $L_0(\alpha)$ at which it first turns negative: $237$ for $\alpha=1$, $897$ for $1.4$, $4590$ for $1.7$, $52\,841$ for $1.9$, $237\,405$ for $1.95$ and $7.35\times10^6$ for $1.99$. The convergence $\alpha_L\to2$ is proved, and it is slow.

\Cref{tab:rate} gives the width rate $c_L(\alpha)$ of \cref{cor:width} and, in parentheses, the width at which the bound $2e^{-c_L(\alpha)n}$ drops below $0.01$. Once the depth is past $L_0(\alpha)$ the rate grows linearly in $L$, and a width of one or two already makes the bound tiny; at $\alpha=1$, $L=300$ the bound is $5.6\times10^{-9}$ at $n=10$ and $3.4\times10^{-43}$ at $n=50$, at $L=1000$ it is $3.8\times10^{-109}$ at $n=10$. The finite-width statement is therefore not weak once it applies; what limits it is the depth $L_0(\alpha)$, not the width.

\begin{table}[h]\centering
\begin{tabular}{rrrrrrr}
\toprule
$\alpha\backslash L$ & 100 & 300 & 1000 & 3000 & $10^4$ & $10^5$\\
\midrule
0.5 & 3.63 (2) & 24.5 (1) & 100.4 (1) & 319.5 (1) & 1089 (1) & 10996 (1)\\
1.0 & -- & 1.97 (3) & 25.0 (1) & 93.3 (1) & 334.8 (1) & 3453 (1)\\
1.5 & -- & -- & -- & 10.6 (1) & 59.2 (1) & 696.5 (1)\\
1.8 & -- & -- & -- & -- & -- & 90.5 (1)\\
1.9 & -- & -- & -- & -- & -- & 11.4 (1)\\
\bottomrule
\end{tabular}
\caption{The width rate $c_L(\alpha)$, with the width at which $2e^{-c_L(\alpha)n}\le0.01$ in parentheses; a dash means the budget is not negative at that depth.}
\label{tab:rate}
\end{table}

On the other side, \cref{thm:above2} gives $\PP(K_{n,L}(\alpha)\le1)\le5\alpha^2L/(n(\alpha-2)^2)$: at $\alpha=3$ and $L=10$ this is $0.45$ at $n=1000$ and $0.045$ at $n=10^4$; at $\alpha=4$, $L=10$, it is $0.2$ and $0.02$. The Chebyshev bound decays like $1/n$; the Chernoff form would decay exponentially.

\subsection{What small networks do}

\Cref{tab:small} reports the measured Lipschitz constant $K_{n,L}(1)$ at small width. By the scaling identity \cref{eq:scaling}, $K_{n,L}(\alpha)\le1$ exactly when $\alpha\le K_{n,L}(1)^{-2/L}$, so the median of $K_{n,L}(1)^{-2/L}$ is the variance parameter at which half the networks are non-expansive. At $n=2$ the constant is exact, by sweeping all directions of the plane; for $n\ge3$ it is a sampled lower bound, so the threshold column is an upper bound. At depth one the operator norm of $W_1$ approaches $2$ at $\alpha=1$ (measured $1.995$ at $n=2000$), so the threshold approaches $\tfrac14$.

\begin{table}[h]\centering
\begin{tabular}{rrrrrl}
\toprule
$n$ & $L$ & median $K_{n,L}(1)$ & 10\%--90\% & median threshold $K^{-2/L}$ & method\\
\midrule
2 & 1 & 1.253 & 0.72--1.84 & 0.64 & exact\\
2 & 2 & 0.851 & 0.25--1.78 & 1.18 & exact\\
2 & 3 & 0.336 & 0.00--1.52 & 2.07 & exact\\
2 & 4 & 0.065 & 0.00--1.02 & 3.93 & exact\\
3 & 2 & 1.336 & 0.65--2.36 & 0.75 & sampled\\
3 & 3 & 0.876 & 0.24--1.91 & 1.09 & sampled\\
5 & 3 & 1.442 & 0.97--2.25 & 0.78 & sampled\\
8 & 3 & 2.067 & 1.50--2.81 & 0.62 & sampled\\
\bottomrule
\end{tabular}
\caption{The global Lipschitz constant at $\alpha=1$ for small networks (400 draws at $n=2$, fewer at larger $n$).}
\label{tab:small}
\end{table}

Two features of \cref{tab:small} are worth reading against the theorem. At width two, deeper networks die: with probability growing in $L$ some layer maps everything to zero, $K=0$, and the median threshold rises above $2$ for that reason alone, which is a small-width effect the theorem says nothing about. At width eight and depth three the median threshold is about $0.6$, well below $2$, in keeping with the depth $L_0$ the budget demands before the bound takes effect. The theorem is about the joint regime of large width first and large depth after; the table is a reminder of what it does not say.

\subsection{One matrix, many steps}

\Cref{tab:horizon} measures the statistic of \cref{thm:fh}, the worst relative deviation $\max_{t\le8}\abs{2^t\norm{x_t}^2-1}$ over eight steps, for the tied dynamics of \cref{sec:horizon} and, beside it, for the feedforward pass with a fresh matrix at every step (the object of \cref{sec:intro} at $\alpha=1$); 200 draws per width. Both columns shrink with the width, as both theorems say. The tied matrix wanders more than fresh matrices at this horizon: at $N=2000$ six percent of the tied trajectories still stray by more than one half, against half a percent for fresh matrices. The theorem's rate is exponential in $N$ with constants that depend on $T$, and eight steps of one matrix are not eight independent draws.

\begin{table}[h]\centering
\begin{tabular}{rrrrrrr}
\toprule
 & \multicolumn{3}{c}{one matrix} & \multicolumn{3}{c}{fresh matrix per step}\\
$N$ & median & 90\% & frac.\ $>\tfrac12$ & median & 90\% & frac.\ $>\tfrac12$\\
\midrule
20 & 0.983 & 12.93 & 0.935 & 0.917 & 2.663 & 0.965\\
50 & 0.861 & 4.734 & 0.865 & 0.749 & 1.903 & 0.835\\
100 & 0.719 & 2.556 & 0.685 & 0.512 & 1.159 & 0.545\\
500 & 0.348 & 0.789 & 0.270 & 0.274 & 0.547 & 0.125\\
2000 & 0.192 & 0.434 & 0.060 & 0.137 & 0.243 & 0.005\\
\bottomrule
\end{tabular}
\caption{The worst deviation $\max_{t\le8}|2^t\|x_t\|^2-1|$ over eight steps, 200 draws per width, seed fixed.}
\label{tab:horizon}
\end{table}

\Cref{tab:chernoff} evaluates the rate $I(1/\alpha)$ of \cref{cor:chernoff} by a grid search over $s$, and compares the two bounds on $\PP(K_{n,L}(\alpha)\le1)$ at $L=10$, $n=1000$. Near $\alpha=2$ both are useless at this width; from $\alpha=3$ on the exponential bound is smaller by five, fifteen and thirty orders of magnitude.

\begin{table}[h]\centering
\begin{tabular}{rrrr}
\toprule
$\alpha$ & $I(1/\alpha)$ & Chebyshev, $5\alpha^2L/(n(\alpha-2)^2)$ & Chernoff, $Le^{-nI(1/\alpha)}$\\
\midrule
2.2 & 0.00087 & 6.05 & 4.2\\
2.5 & 0.00455 & 1.25 & 0.11\\
3.0 & 0.01395 & 0.45 & $8.7\times10^{-6}$\\
4.0 & 0.03642 & 0.20 & $1.5\times10^{-15}$\\
6.0 & 0.07847 & 0.11 & $8.4\times10^{-34}$\\
\bottomrule
\end{tabular}
\caption{The expansive side at $L=10$, $n=1000$: the Chebyshev bound of \cref{thm:above2} against the exponential bound of \cref{cor:chernoff}.}
\label{tab:chernoff}
\end{table}

\section{Sharpness, and some structural remarks}
\label{sec:sharp}

The finite theorem uses the elementary function $g_\alpha$ of \cref{lem:scalar}. This section shows that the same function is the exact limiting moment exponent, so the scalar step loses no exponential rate at large width. It then records some structural consequences of the gauge identity.

\subsection{The exact scalar high-moment exponent}

For $\theta>0$ and $0\le u\le1$ define
\begin{equation}\label{eq:Phi}
\Phi_{\alpha,\theta}(u)=H(u)-\log2+\theta\log(2\alpha)+\Bigl(\frac u2+\theta\Bigr)\log\Bigl(\frac u2+\theta\Bigr)-\frac u2\log\frac u2-\theta,
\end{equation}
continuous on $[0,1]$ because $x\log x$ extends continuously to zero.

For a chi-square variable with $m\ge1$ degrees of freedom and $r>0$, the substitution $x=2y$ in $\int_0^\infty x^r\frac{x^{m/2-1}e^{-x/2}}{2^{m/2}\Gamma(m/2)}\dd x$ gives
\begin{equation}\label{eq:chi-moment}
\E[(\chi^2_m)^r]=2^r\,\frac{\Gamma(m/2+r)}{\Gamma(m/2)} .
\end{equation}
Conditioning on $M_n$ in \cref{eq:mixture} with $r=\theta n$ (the term $m=0$ contributes zero),
\begin{equation}\label{eq:exact-sum}
\E[(X_n^{(\alpha)})^{\theta n}]=2^{-n}\sum_{m=1}^n\binom nm\Bigl(\frac{2\alpha}n\Bigr)^{\theta n}\frac{\Gamma(m/2+\theta n)}{\Gamma(m/2)} .
\end{equation}
Two forms of Stirling's formula are used: for integers $k\ge1$, $\log k!=(k+\tfrac12)\log k-k+O(1)$, and for real $x\ge\tfrac12$, $\log\Gamma(x)=(x-\tfrac12)\log x-x+O(1)$ uniformly in $x\ge\tfrac12$ (the function $\log\Gamma(x)-(x-\tfrac12)\log x+x$ is continuous on $[\tfrac12,\infty)$ and tends to $\tfrac12\log2\pi$). With $u=m/n$ they give, uniformly for $1\le m\le n$,
\begin{gather*}
\log\binom nm=nH(u)+O(\log n),\\
\log\Gamma\Bigl(n\bigl(\tfrac u2+\theta\bigr)\Bigr)-\log\Gamma\Bigl(\frac{nu}2\Bigr)=\theta n\log n+n\Bigl[\bigl(\tfrac u2+\theta\bigr)\log\bigl(\tfrac u2+\theta\bigr)-\tfrac u2\log\tfrac u2-\theta\Bigr]+O_\theta(\log n).
\end{gather*}
The term $-\theta n\log n$ from $(2\alpha/n)^{\theta n}$ cancels the $+\theta n\log n$ of the Gamma ratio, and the $m$-th summand $a_{n,m}$ of \cref{eq:exact-sum} satisfies $\frac1n\log a_{n,m}=\Phi_{\alpha,\theta}(m/n)+O_{\alpha,\theta}(\log n/n)$ uniformly. A sum of $n$ nonnegative terms lies between its largest term and $n$ times it, so after taking $\frac1n\log$ the sum and the maximum differ by at most $\log n/n$; and the grids $\{1/n,\dots,1\}$ become dense, so the grid maximum of the continuous $\Phi_{\alpha,\theta}$ converges to its maximum.

\begin{proposition}[exact exponent]\label{prop:exact}
For every $\alpha>0$ and $\theta>0$,
\[
\lim_{n\to\infty}\frac1n\log\E[(X_n^{(\alpha)})^{\theta n}]
=\max_{0\le u\le1}\Phi_{\alpha,\theta}(u)=g_\alpha(\theta),
\]
where $g_\alpha$ is the finite-width bound in \cref{eq:scalar-def}. The unique maximizing fraction is $u_\theta=1/(1+z_\theta)$.
\end{proposition}

\begin{proof}
The limit and its variational formula follow from the exact Gamma sum and uniform estimates above. For $0<u<1$,
\[
\partial_u\Phi=\log\frac{1-u}{u}+\frac12\log\frac{u/2+\theta}{u/2},\qquad
\partial_{uu}\Phi=-\frac1{u(1-u)}+\frac1{4(u/2+\theta)}-\frac1{2u}<0.
\]
The first derivative tends to $+\infty$ at zero and $-\infty$ at one, so there is one interior maximizer. Substitute $u=1/(1+z)$. Its stationary equation reduces to $\theta=(1-z)/(2z^2)$, giving $z=z_\theta$ and
\[
\frac{u_\theta}2+\theta=\frac1{2z_\theta^2(1+z_\theta)}.
\]
Substituting these identities into \cref{eq:Phi} and collecting logarithms gives exactly \cref{eq:scalar-def}. This identifies the limiting exponent with the function that already bounds every finite-width moment in \cref{lem:scalar}.
\end{proof}

In particular $u_\theta=\tfrac12+\tfrac12\theta+O(\theta^2)$ and \cref{eq:scalar-expansion} gives
\begin{equation}\label{eq:g-expansion}
g_{2e^{-\delta}}(\theta)=-\delta\theta+\frac52\theta^2+O(\theta^3).
\end{equation}
The scalar exponent is exact, but the geometric count and the sphere net in the master inequality still introduce slack. The constant $\sqrt{10}$ remains a constant in a certified lower bound for the network threshold, not a matching asymptotic description of that threshold.

\subsection{The Cover factor at fixed depth}

\begin{lemma}\label{lem:C-rate}
For every fixed $L\ge2$, $\lim_{n\to\infty}\frac1n\log C_{n,L}=LH(1/L)$.
\end{lemma}

\begin{proof}
For $L>2$ the ratio of consecutive terms $\binom{nL-1}{j+1}/\binom{nL-1}{j}=(nL-1-j)/(j+1)$ exceeds $1$ for $0\le j\le n-2$, because $nL-1-j\ge nL-n+1>n-1\ge j+1$; so the last term is the largest and $\binom{nL-1}{n-1}\le C_{n,L}/2\le n\binom{nL-1}{n-1}$. Stirling's formula gives $\frac1n\log\binom{nL-1}{n-1}\to LH(1/L)$, and the factors $2$ and $n$ vanish after $\frac1n\log$. For $L=2$, $C_{n,2}=2^{2n-1}$ and $\frac1n\log C_{n,2}\to2\log2=2H(\tfrac12)$.
\end{proof}

So the entropy bound of \cref{lem:C-entropy} is exact at the exponential scale, and the finite theorem loses nothing there at large width. At small width it does lose something; \cref{sec:numerics} shows how much.

\subsection{The all-active pattern, exactly}

Take $D_1=\dots=D_L=I_n$. Then $H^D_\ell=Q_\ell:=W_\ell\cdots W_1$ and $A_D$ stacks $Q_1,\dots,Q_L$. Every Gaussian square matrix is invertible almost surely, since its determinant is a nonzero polynomial in continuously distributed entries and the zero set of a nonzero polynomial has Lebesgue measure zero (\cref{lem:poly}). On that event $(W_1,\dots,W_L)\mapsto(Q_1,\dots,Q_L)$ is a smooth bijection with inverse $W_1=Q_1$, $W_\ell=Q_\ell Q_{\ell-1}^{-1}$, so the tuple $(Q_\ell)$ has a joint density on $(GL_n)^L$. Choose any $n$ of the $nL$ rows of the stacked matrix; their determinant is a polynomial in the entries of the $Q_\ell$, not identically zero (assign the chosen rows the standard basis vectors and complete each $Q_\ell$ to an invertible matrix), hence nonzero almost surely; finitely many choices, so all $n\times n$ minors are nonzero almost surely. That is central general position for the coordinate hyperplanes restricted to $\im A_D$, and the region count is then exact: $r(\im A_D)=C_{n,L}$ almost surely. \Cref{prop:orbit} with $F=1$ gives
\begin{equation}\label{eq:all-active}
\PP(R_D\ne\varnothing)=\frac{C_{n,L}}{2^{nL}}=2^{1-nL}\sum_{j=0}^{n-1}\binom{nL-1}j\qquad(D\text{ all-active}),
\end{equation}
an exact formula for the probability that a random deep network has an input at which every unit is active at every layer. Its logarithm is $-[nL\log2-\log C_{n,L}]=-n[L\log2-LH(1/L)]+o(n)$ by \cref{lem:C-rate}: the accumulated cost per unit width is $L\log2-O(\log L)$, not a constant per layer. For $L=1$ the cost is zero, as it must be: $C_{n,1}=2^n$, and a full-rank $W_1$ reaches every orthant. \Cref{sec:numerics} checks \cref{eq:all-active} against direct enumeration.

\begin{lemma}\label{lem:poly}
If $P:\R^d\to\R$ is a polynomial that is not identically zero, then $\{x:P(x)=0\}$ has $d$-dimensional Lebesgue measure zero.
\end{lemma}

\begin{proof}
Induction on $d$. For $d=1$ a nonzero polynomial has finitely many roots. For $d>1$ write $P=\sum_{k=0}^ma_k(x_1,\dots,x_{d-1})x_d^k$ with some $a_k$ not identically zero. By induction the set where all $a_k$ vanish has measure zero; for every other value of the first $d-1$ coordinates the polynomial in $x_d$ is nonzero and has finitely many roots, a one-dimensional null set. Fubini's theorem gives total measure zero.
\end{proof}

\subsection{Depth asymptotics at fixed width}
\label{sec:fixed-width-depth}

The same identity also controls the opposite order of limits: keep the width fixed and let the depth grow. Write $R_{n,L}=\#\{D:R_D\ne\varnothing\}$ for the strict-region count and $S_{n,L}=\{F^{(\alpha)}_{n,L}\not\equiv0\}$ for global survival. A strictly negative final preactivation can leave a nonempty strict region whose output is zero, so these are different objects.

\begin{corollary}[fixed-width depth asymptotics]\label{cor:fixed-width-depth}
Fix $n\ge1$ and $\alpha>0$, with independent Gaussian layers, no biases, input domain $\R^n$ and ReLU at the output as in \cref{eq:forward}. Put $\beta_n=1-2^{-n}$. There are finite positive constants $\kappa_n$ and $\sigma_n$, independent of $\alpha$, such that
\begin{equation}\label{eq:depth-ratios}
\E R_{n,L}\sim\kappa_n\beta_n^L,
\qquad \PP(S_{n,L})\sim\sigma_n\beta_n^L
\qquad (L\to\infty).
\end{equation}
Both ratios on the left after division by $\beta_n^L$ are nondecreasing. Moreover $\kappa_n\ge n\beta_n^{-2}$ and $\beta_n^{-1}\le\sigma_n\le\kappa_n$. In particular,
\begin{equation}\label{eq:depth-rate}
\lim_{L\to\infty}\frac{\log\E R_{n,L}}L
=\lim_{L\to\infty}\frac{\log\PP(S_{n,L})}L
=\log(1-2^{-n}).
\end{equation}
\end{corollary}

\begin{proof}
Write $q=2^n$, $m=q-1$, $a=q-n-1$ and $\beta=m/q$. For every fixed sequence of nonzero intermediate masks, every row of every formal preactivation matrix is nonzero almost surely. Indeed, conditional on a nonzero retained product $B$, a fresh Gaussian row $w$ satisfies $w^TB\ne0$ almost surely: its zero set is a proper linear subspace. A nonzero mask retains at least one such row. Induction, followed by a finite union over rows and masks, proves the assertion at each depth. These are formal masks; no conditioning on realizability is involved.

Let $A_k$ be the scalar sum of $\E r(\im A_D)$ for the prefix stack $H^D_1,\ldots,H^D_k$, over all $k-1$ intermediate masks of rank at least two. The final mask is omitted because it does not enter the stack. There are $a^{k-1}$ prefixes, and the invertibility of $W_1$ gives
\[
A_1=q,\qquad 0\le A_k\le a^{k-1}C_{n,k}.
\]
At $n=1$, use $0^0=1$: only the empty prefix exists, and $A_k=0$ for $k>1$.

Partition the nonzero intermediate-mask sequences by their first rank-one mask, at layer $k<L$. The selected row of $H^D_k$ already occurs in the prefix stack. Every later row is a nonzero multiple of it almost surely, so later layers add no new hyperplane and leave the orthant count unchanged. There are $n$ choices for the selected row, $m^{L-k-1}$ choices for the later intermediate masks, and $q$ choices for the final mask. This family contributes $nq m^{L-k-1}A_k$ to the formal orthant sum. The family with no rank-one intermediate mask contributes $qA_L$; a zero intermediate mask contributes zero because the next preactivation block vanishes. Summing \cref{eq:orbit} with functional $1$ therefore gives the exact identity
\begin{equation}\label{eq:depth-prefix}
\E R_{n,L}=q^{1-L}\left[A_L+n\sum_{k=1}^{L-1}m^{L-k-1}A_k\right],
\qquad L\ge1.
\end{equation}
After division by $\beta^L$ this is
\[
\frac{\E R_{n,L}}{\beta^L}
=\frac{qA_L}{m^L}+\frac{nq}{m^2}\sum_{k=1}^{L-1}\frac{A_k}{m^{k-1}}.
\]
For $n\ge2$, $\rho=a/m<1$, and $A_k/m^{k-1}\le C_{n,k}\rho^{k-1}$. At fixed $n$, $C_{n,k}$ is a polynomial in $k$ of degree $n-1$. The first term tends to zero and the series converges, proving
\begin{equation}\label{eq:depth-constant}
\kappa_n=\frac{nq}{m^2}\sum_{k\ge1}\frac{A_k}{m^{k-1}}\in(0,\infty).
\end{equation}
The first summand gives $\kappa_n\ge nq^2/m^2=n\beta^{-2}$. The same formula gives $\kappa_1=4$ directly. Extending any prefix by one of the $a$ admissible masks and a fresh layer only appends proper hyperplanes, so it cannot remove an old open region. Thus $A_{k+1}\ge aA_k$, and subtraction in \cref{eq:depth-prefix} yields
\[
\E R_{n,L+1}-\beta\E R_{n,L}
=q^{-L}(A_{L+1}-aA_L)\ge0.
\]
This proves monotonicity of the normalized means.

For survival, couple the depths by one infinite sequence of independent matrices and put $p_L=\PP(S_{n,L})$. On $S_{n,L}$, choose the first input in a fixed enumeration of $\mathbb Q^n$ at which $F^{(\alpha)}_{n,L}$ is nonzero. Such an input exists by continuity, and the choice is measurable with respect to $W_1,\ldots,W_L$ because the enumeration is countable. Conditional on this past, the next fresh Gaussian layer preserves its nonzero output with probability exactly $\beta$. Its survival implies global survival, so $p_{L+1}\ge\beta p_L$; equality is not asserted. Positive homogeneity gives $F(0)=0$, and \cref{lem:reduction} ensures that a nonzero network has a realized strict region. Consequently $p_L\le\E R_{n,L}$. The sequence $p_L/\beta^L$ is therefore nondecreasing and bounded above by $\kappa_n$. Since $p_1=1$ almost surely, it converges to $\sigma_n\in[\beta^{-1},\kappa_n]$. Finally, a positive common rescaling of Gaussian weights changes neither strict signs nor zero-output events, so both constants are independent of $\alpha>0$.
\end{proof}

The finite bounds behind these asymptotics are also explicit:
\begin{align}
n\beta_n^{L-2}\ \le\ \E R_{n,L}\ &\le\ C_{n,L}\beta_n^{L-1}, &&L\ge2,\label{eq:depth-strict}\\
\beta_n^{L-1}\ \le\ \PP(S_{n,L})\ &\le\ \min\{1,C_{n,L}\beta_n^L\}, &&L\ge1.\label{eq:depth-survival}
\end{align}
For the mean, the lower bound is the $k=1$ term in \cref{eq:depth-prefix}; the upper bound counts the $q m^{L-1}$ masks with nonzero intermediate masks and bounds each orthant count by $C_{n,L}$. For global survival, use \cref{prop:orbit} with the invariant functional $\one_{\{J_D\ne0\}}$: a zero mask at any layer, including the final one, kills the Jacobian, leaving at most $m^L$ contributors. For the lower bound choose $s=W_1^{-1}e_1$, so $x_1(s)=e_1$, and propagate this first-layer witness through the independent remaining layers. Its survival probability is $\beta_n^{L-1}$. A deterministic nonzero input instead survives with probability exactly $\beta_n^L$.

At width one, direct sign enumeration gives $\E R_{1,L}=2^{2-L}$ and $\PP(S_{1,L})=2^{1-L}$, hence $(\kappa_1,\sigma_1)=(4,2)$. At $L=2$ there is one strict region for every nonzero scalar pair, while global survival has probability $1/2$. The constants at higher widths are left implicit, except for $\kappa_2=8$ below. These are fixed-width limits; no uniformity in a growing width or quantitative convergence rate for the survival ratio is asserted.

Earlier work already establishes death at large depth. In the zero-bias square setting, Rister and Rubin~\cite{rister2021}, Sections~3--4, give a global survival upper bound with base $1-2^{-n^2}$ and a fixed-input lower bound with base $\beta_n$. Their Section~5.3 also studies a survival ratio for finite data as width grows, under a variance assumption. Lu et al.~\cite{lu2020}, Theorem~3.5 and Appendix~C, give two-exponential bounds and a one-ray state analysis for one-dimensional input, with a final affine readout. After accounting for that depth convention, their bounds and the past-measurable witness argument above also imply a finite positive survival-ratio limit in the scalar-input setting. For any fixed finite set of $M$ nonzero inputs, the simpler bound $p_L\le M\beta_n^L$ and the same monotonicity argument suffice. The additional geometric step here bounds the normalized mean strict count and hence survival on the entire square input space $\R^n$. Neither qualitative eventual death nor ratio arguments in general are claimed as new.

\subsection{The mean strict-region count at width two}

Let $R_{2,L}$ be the number of nonempty strict regions $R_D$ in a width-two network. This counts activation patterns with every preactivation nonzero; it excludes ordinary cells on which a preactivation is identically zero.

\begin{proposition}\label{prop:width-two-mean}
For every $L\ge1$ and every $\alpha>0$,
\begin{equation}\label{eq:width-two-mean}
\E R_{2,L}=8\left[\left(\frac34\right)^L-\left(\frac14\right)^L\right].
\end{equation}
In particular, $\kappa_2=8$ in \cref{eq:depth-ratios}.
\end{proposition}

\begin{proof}
At width two, the only intermediate mask of rank at least two is $I_2$. Thus the prefix in \cref{eq:depth-prefix} stacks the unmasked products $Q_1,\ldots,Q_k$. Their $2k$ row directions are distinct almost surely by the joint-density argument for \cref{eq:all-active}. They define $2k$ lines in the input plane, so $A_k=4k$. Substituting $q=4$, $m=3$ and $n=2$ in \cref{eq:depth-prefix} gives
\[
\E R_{2,L}=4^{-L}\left[32\sum_{k=1}^{L-1}k3^{L-k-1}+16L\right]
=8\frac{3^L-1}{4^L}.
\]
The finite-sum identity follows by differentiating a geometric sum, including $L=1$ when the sum is empty. Dividing by $(3/4)^L$ gives the limit $8$.
\end{proof}

The first six means are $4,4,3.25,2.5,1.890625,1.421875$. Their decay at large depth concerns strict regions. Ordinary cells remain even when the network is identically zero, so \cref{eq:width-two-mean} must not be read as their mean count. No affine-bias version of this formula is asserted.

\subsection{Depth one}

At $L=1$ the all-active mask is realized with probability one, $J_D=W_1$ for it, and \cref{eq:reduction} gives $K_{n,1}(\alpha)=\op{W_1}$: every mask $D_1$ gives $\op{D_1W_1}\le\op{W_1}$. By the Bai--Yin theorem \cite{baiyin}, for the centered Gaussian matrix of this paper, with independent $\mathcal N(0,\alpha/n)$ entries, $\op{W_1}\to2\sqrt\alpha$ almost surely. Hence $\PP(K_{n,1}(\alpha)\le1)\to1$ for $\alpha<\tfrac14$ and $\to0$ for $\alpha>\tfrac14$, that is,
\[
\alpha_1=\tfrac14 .
\]
A depth-one network at $\alpha=1$ is expansive with probability tending to one, although $1<2$; there is no contradiction with \cref{thm:main}, whose part (i) only takes effect beyond $L_0(\alpha)$.

\subsection{Why the global upper bound never conditions on an input}

A tempting route would be to choose an input, condition on the activation masks it generates, and analyse the resulting Jacobian as if the masks were independent of the weights. The last step is false in general: the event that a mask is generated by an input carries information about the same weights that appear in the Jacobian. The gauge argument fixes a formal mask without choosing a witness input, multiplies the realizability indicator by a gauge-invariant Jacobian functional, and averages over a finite symmetry group exactly; only then does it sum over masks. The sphere net is used for tangent directions of a fixed linear map, not to search for regions, so a realizable cone of arbitrarily small angle is still counted.

\subsection{The order of limits}

The master inequality \cref{eq:master} and the finite theorem are statements about every $n$ and $L$, and one may evaluate them in any joint regime. What is proved about the thresholds, however, fixes $L$ first. If $L=L(n)$ grows rapidly, several effects change scale; for instance the probability that some layer dies at a fixed input is at most $L2^{-n}$, negligible for fixed $L$ but not for every $L(n)$. No joint-limit statement should be read into \cref{thm:main} beyond what \cref{thm:finite} gives when its budget is evaluated in that regime.


\section*{Acknowledgments and funding}
The research presented in this paper was supported by the European Research Council (ERC) under the European Union's Horizon Europe research and innovation programme (grant agreement No.~101041711), by the Simons Foundation as part of the Collaboration on the Mathematical and Scientific Foundations of Deep Learning, by Heights Labs, by Convex Nexus Capital, by the Israel Science Foundation (grant number 2258/19), and by the Israel Science Foundation (ISF Grant 4101/25).

\section*{Declaration of competing interest}
The author holds equity in Heights Labs and in Convex Nexus Capital and has a management role at Convex Nexus Capital; both entities supported this work and are named above. Neither had a role in the design of the study, the analysis, the writing, or the decision to publish. The public funders had no role in those either.

\section*{Declaration of generative AI and AI-assisted technologies in the manuscript preparation process}
During the preparation of this work the author made substantial use of generative AI tools, and describes their role here. The two theorems and their proofs are the author's, from two earlier manuscripts. The present text was rewritten and consolidated with Fable~5.1 (Anthropic, through the Claude Code agent, September 2026): the finite-width formulation of the first theorem and its computable threshold, the condensed account of the finite-horizon proof, the section comparing the two arguments and the corollary that sharpens the expansive side, the extension to biases, the numerical experiments and the code behind every table, and the drafting of the exposition. Subsequent Codex-assisted revisions corrected the strict-region enumeration, checked the affected numerical calculations, distinguished sufficient certificates from necessary conditions, added a comparison with related work, and repaired the bias hypotheses and proofs. The released source and code record these changes and their verification scope. The author takes responsibility for the manuscript; computational checks do not replace mathematical review.

\section*{Data and code availability}
The code behind every table, the fixed-seed simulation outputs, and the scripts that evaluate the exponent budget, the thresholds $\alpha^*_L$ and the rates $I(1/\alpha)$ are included in the source archive and available at \url{https://github.com/yspennstate/relu-global-stability-numerics}.

\begin{thebibliography}{99}
\bibitem{geuchen2025} P. Geuchen, D. St\"oger, T. Telaar and F. Voigtlaender, Upper and lower bounds for the Lipschitz constant of random neural networks, \href{https://arxiv.org/abs/2311.01356v4}{arXiv:2311.01356v4}, 2025.
\bibitem{dirksen2025} S. Dirksen, P. Finke, P. Geuchen, D. St\"oger and F. Voigtlaender, Near-optimal estimates for the $\ell^p$-Lipschitz constants of deep random ReLU neural networks, \href{https://arxiv.org/abs/2506.19695v1}{arXiv:2506.19695v1}, 2025.
\bibitem{yang2019} G. Yang, Tensor Programs I: Wide Feedforward or Recurrent Neural Networks of Any Architecture are Gaussian Processes, \emph{NeurIPS}, 2019; \href{https://arxiv.org/abs/1910.12478v3}{arXiv:1910.12478v3}, 2021.
\bibitem{hanin2019} B. Hanin and D. Rolnick, Complexity of Linear Regions in Deep Networks, \emph{ICML}, 2019; \href{https://arxiv.org/abs/1901.09021v2}{arXiv:1901.09021v2}.
\bibitem{wang2022} Y. Wang, Estimation and Comparison of Linear Regions for ReLU Networks, \emph{IJCAI}, 2022, 3544--3550; \href{https://www.ijcai.org/proceedings/2022/0492.pdf}{proceedings paper}.
\bibitem{lu2020} L. Lu, Y. Shin, Y. Su and G.~E. Karniadakis, Dying ReLU and Initialization: Theory and Numerical Examples, \href{https://arxiv.org/abs/1903.06733v3}{arXiv:1903.06733v3}, 2020.
\bibitem{cover} T.~M. Cover, Geometrical and statistical properties of systems of linear inequalities with applications in pattern recognition, \emph{IEEE Trans. Electronic Computers} EC-14 (1965), 326--334.
\bibitem{schlafli} L. Schl\"afli, \emph{Theorie der vielfachen Kontinuit\"at}, 1852 (published 1901); the region count for central hyperplane arrangements.
\bibitem{vershynin} R. Vershynin, \emph{High-Dimensional Probability}, Cambridge University Press, 2018; sphere nets and operator norms.
\bibitem{cover-thomas} T.~M. Cover and J.~A. Thomas, \emph{Elements of Information Theory}, 2nd ed., Wiley, 2006; the entropy bound for binomial sums.
\bibitem{baiyin} Z.~D. Bai and Y.~Q. Yin, Necessary and sufficient conditions for almost sure convergence of the largest eigenvalue of a Wigner matrix, \emph{Ann. Probab.} 16 (1988), 1729--1741; and Limit of the smallest eigenvalue of a large dimensional sample covariance matrix, \emph{Ann. Probab.} 21 (1993), 1275--1294.
\end{thebibliography}

\end{document}
